\subsection{The Albanese coordinates and real structures}
\begin{proposition}[exact nomes and continuation radius]\label{ari:C.2.1}
\StatementRecord{ari:C.2.1}{proved}{Historical events: Claude Fable 5.2 (Anthropic), effort high, 16 November 2026: used and rechecked in this run. Claude Fable 5.2 (Anthropic), effort high, 17 November 2026: carried as background without fault found. GPT-6 Astra (OpenAI), ChatGPT Pro mode, effort Pro, S6\_09\_PAPERS, 17 November 2026: targeted proof check. Also checked in this writing. Pre-continuation writing events: GPT-6 Astra (OpenAI), ChatGPT Pro mode, effort Pro, S6\_10\_PAPERS, 18 November 2026: computation or formal certificate; targeted proof check. Continuation mathematical and source checks: GPT-6 Astra (OpenAI), ChatGPT Pro mode, effort Pro, S6\_10\_PAPERS continuation, 19 November 2026: computation or formal certificate; independent re-proof; used and rechecked in this run. Continuation metadata events: GPT-6 Astra (OpenAI), ChatGPT Pro mode, effort Pro, S6\_10\_PAPERS continuation, 19 November 2026: used and rechecked in this run; status-vocabulary reconciliation; metadata reconciliation only; no new mathematical check. \emph{Continuation recorded proof scope.} Complete new mathematical derivation written before opening the target proof, followed by comparison of every new step. The fresh program \path{code/continuation_root/ari_C2_C3_fresh.py} passed 34 exact identities by a differential-polynomial and real-polynomial method. Ramanujan, theta and period inputs were opened with their exact normalizations; the renewed credit search and positive antecedents are recorded in the continuation comparison.}{the hypotheses, notation and definitions stated here; \cref{fam:src:A.4.1}; \cref{ari:C.1.2}; \cref{fam:src:C.2.3,fam:src:A.5.1}; \cite{DuGuoKusakabeXie2026v1}}

For $w=e^{-\pi ic}$, the fixed branches give
\begin{align*}
 \kappa_1&=c+\tfrac12-\frac{3i}{2\pi}\log48,&
 q_1&=-w^2/48^3,\\
 \kappa_2&=c+\tfrac32-\frac{3i}{\pi}\log6,&
 q_2&=iw/216,\qquad q_1=(27/64)q_2^2.
\end{align*}
The family is parameterized by $0<|q_2|<1$ and exists precisely in
that domain. Writing $J(q)=q^{-1}+744+\cdots$, one has
\[
 j(\operatorname{Alb}S_1)=J(-w^2/48^3).
\]
This function is holomorphic on $0<|w|<48^{3/2}$, has a pole of order
two at zero, and has a natural boundary at $|w|=48^{3/2}$.
In particular it continues holomorphically through every point of the
family boundary $|w|=216$.
\end{proposition}


\begin{proof}
\begin{enumerate}[label=\arabic*.,leftmargin=*]
\item The precise Albanese lattices in \cref{fam:src:A.4.1} are verified by projecting all
translations on the product covers. At the order-three fibre the
projection is $(1/3)(\mathbb Z+\kappa_1\mathbb Z)$; multiplication by
three gives the first stated lattice. At the order-four fibre it is
$\tfrac12\mathbb Z+\tfrac14\kappa_2\mathbb Z$; multiplication by two
gives $\mathbb Z+(\kappa_2/2)\mathbb Z$. Thus the positive-imaginary
period representatives are $-\kappa_1$ and $-\kappa_2/2$.
Adding $2\mu_1=1-i/\sqrt3$ and
$3\mu_2=3/2-3i/2$ to \ref{ari:C.1.2} proves the two $\kappa$ identities,
including their real parts. Exponentiation gives the signs in $q_1,q_2$;
the ratio is $216^2/48^3=27/64$.

\item The identity $q_2=iw/216$ in Proposition~\ref{push:src:D.1.1} has
factor $1/2$ precisely because of the Albanese projection just computed.
Its natural-boundary proof uses modular invariance at rational boundary
points. The existence criterion and exact $2\mathbb Z$
identification are \cref{fam:src:C.2.3,fam:src:A.5.1}; the latter is also Theorem~1.4
(label \texttt{p2:thm:exact-period}) of
\cite{DuGuoKusakabeXie2026v1}. Its hypotheses specify precisely the
retained period family, affine finite twists and cusp collar. They
coincide with the marked family $X_c$; the normalization differs
only by the fixed choice of the particular solution, which does not
alter parameter differences. Version 2 drops that theorem; we infer
no mathematical conclusion from this omission.

\item The Fourier expansion for $J$ converges for $0<|q|<1$ and has a simple
pole at $q=0$. The inverse image under $q=-w^2/48^3$ is the claimed
punctured disc and the pole has order two. The inequality
$216<48^{3/2}$ is equivalent to $27<64$. At every nonzero point on the
larger boundary the map to $q$ is locally biholomorphic. A continuation
there would continue $J$ through the unit circle, contradicting the
rational-cusp argument proved in Proposition~\ref{push:src:D.1.1}. The phrase ``up to the
radius'' therefore means the open punctured disc, not holomorphy on
its bounding circle or at its central pole.
\end{enumerate}
\end{proof}

\begin{proposition}[the real locus]\label{ari:C.2.2}
\StatementRecord{ari:C.2.2}{proved}{Historical events: Claude Fable 5.2 (Anthropic), effort high, 15 November 2026: carried as background without fault found. Claude Fable 5.2 (Anthropic), effort high, 16 November 2026: carried as background without fault found; used and rechecked in this run. GPT-6 Astra (OpenAI), ChatGPT Pro mode, effort Pro, S6\_09\_PAPERS, 17 November 2026: targeted proof check. Also checked in this writing. Pre-continuation writing events: GPT-6 Astra (OpenAI), ChatGPT Pro mode, effort Pro, S6\_10\_PAPERS, 18 November 2026: computation or formal certificate; targeted proof check. Continuation mathematical and source checks: GPT-6 Astra (OpenAI), ChatGPT Pro mode, effort Pro, S6\_10\_PAPERS continuation, 19 November 2026: targeted proof check. Continuation metadata events: GPT-6 Astra (OpenAI), ChatGPT Pro mode, effort Pro, S6\_10\_PAPERS continuation, 19 November 2026: used and rechecked in this run; status-vocabulary reconciliation; metadata reconciliation only; no new mathematical check. \emph{Pre-continuation recorded proof scope.} Re-derived the exact reflection intertwiners, primitive norm vectors, Laurent gluing, conjugation action on translation C/Z and the real square correction. The full-automorphism source pointer is being corrected to the credited local proof. Fully read and freshly replayed unchanged Cspecial\_reflection.py; exit zero. Output and SHA256 are retained in code/I/replay\_special/replay\_manifest.json. Numerical outputs retain numerical scope.}{the hypotheses, notation and definitions stated here; \cref{fam:src:B.4.2}; \cref{fam:src:A.5.1}; \cite{DuGuoKusakabeXie2026v1}}

An antiholomorphic equivalence $X_c\to X_{c'}$ exists exactly when
$c'-(1-\bar c)\in2\mathbb Z$. Consequently $X_c$ admits an
antiholomorphic involution exactly when $q_2\in\mathbb R\setminus\{0\}$.
\end{proposition}


\begin{proof}
\begin{enumerate}[label=\arabic*.,leftmargin=*]
\item We check the additional reflection argument, rather than infer a real
structure from the reality of one elliptic invariant. On the period
cover, reflection acts by
\[
 (\tau,\mu,\beta)\longmapsto(-\bar\tau,1-\bar\mu,-\bar\beta).
\]
The theta product and the cusp normalization prove this equality of
continued period functions. For the integral monodromy matrices
$A_1,A_2$ of the retained construction, the lattice comparison is
\[
 H=\begin{pmatrix}1&0&0&0\\0&1&0&0\\-6&0&-1&0\\0&-1&0&-1\end{pmatrix},
 \quad H^2=I,\quad HA_2H=A_2^{-1},\quad
 HA_1H=A_2^{-1}A_1^{-1}A_2,
\]
\[
 \Pi_{-\bar c}(\iota z)H=-\overline{\Pi_c(z)}.
\]
Thus $(z,\zeta)\mapsto(\iota z,-\bar\zeta)$ is an actual torus-family
map. Write $d=(0,0,0,1)^t$, $\varepsilon=(1,2,-4,0)^t$ and
$\varepsilon'=(1,3,-3,0)^t$. Reversing the finite deck generators
changes the target affine vectors to
\[
 v'_1=-\varepsilon+d,\qquad v'_2=\varepsilon'-3d.
\]
The second equality includes the sign of the inverse affine action.
The file \path{Cspecial_reflection.py} checks these matrix identities
over the integers.

\item Normalize this reflected datum to the retained vectors
$(\varepsilon,-\varepsilon')$. The linear fibre map is $-I_2$, the
parameter shift is $1$, and the lattice map is $-I+E$,
$E=d(1,0,0,0)$. The differences between target and transported affine
vectors are $2d$ and $-4d$. They lie in the images of the respective
norm operators. Explicit integral correcting cocycles are
$\ell_1=(0,0,-2,0)^t$ and $\ell_2=(0,2,0,0)^t$; the norm equations
are $N_1\ell_1=-2d$ and $N_2\ell_2=4d$. Rational solutions of
$(I-A_j)r_j=(v_j-(-I+E)v'_j)/m_j+\ell_j$ can be chosen with first
coordinate zero: $r_1=(0,-2/3,-2/3,0)^t$ and
$r_2=(0,1,1,0)^t$. These give translations holomorphic across both
root discs. Their cusp cocycle is
$\ell_0=-A_0(\ell_1+A_1\ell_2)=(0,0,4,-2)^t$; its coordinates are
integral, so the corresponding shear and inversion preserve the
periodic triangular cusp fan. They therefore extend across its toric
charts. The differences of the local corrections form a cocycle in
the rank-two vector bundle whose invariant line and quotient are
$\mathcal O$ and $\mathcal O(-1)$ on the base sphere. Both have zero
first cohomology, so the local maps glue. Equivalently one splits the
Laurent series of the cocycle into its nonnegative and negative
parts in these two line-bundle frames. This checks the precise
gluing step of Proposition~\ref{fam:src:B.4.2}, including the finite affine translations
and cusp extension. The resulting map has target $X_{1-\bar c}$.

\item Every other antiholomorphic comparison composed with its inverse is
holomorphic, so \cref{fam:src:A.5.1} gives exactly the congruence in the statement.
A self-map exists iff $\operatorname{Re}c\in1/2+\mathbb Z$, iff
$w=-\bar w$, iff $q_2=iw/216$ is real. Its square is a holomorphic
translation $d_0\in\mathbb C/\mathbb Z$, since the automorphism group
is the $\mathbb C^*$ computed in the credited re-derivation
of Proposition~\ref{fam:src:B.4.1}, Appendix~\ref{app:rederive:fam-src-B.4.1}. The constructed reflection acts on
this additive translation coordinate by conjugation. Commutation
with its square gives $d_0=\bar d_0$ in $\mathbb C/\mathbb Z$, hence
$d_0$ has a real representative. Composing the reflection with
translation by $-d_0/2$ makes its square the identity. Necessity and
sufficiency for an involution follow.
\end{enumerate}
\end{proof}

\subsection{The next local terms}
Put $G=g_2(i)=\Gamma(1/4)^8/(16\pi^2)$ and
$D_*=-3\log6/\pi$.
Our theta convention is
$\vartheta_j(z,\tau)=\theta_j(\pi z\mid\tau)$, with lattice
$\mathbb Z+\tau\mathbb Z$.  The heat identity and modular
transformations are \cite[Equations 20.13.1, 20.7.26, 20.7.30
and 23.18.5]{DLMFTheta2026}.  The classical differential system,
in the notation $P=E_2$, $Q=E_4$, $R=E_6$, is
\[
 xP_x=(P^2-Q)/12,\qquad
 xQ_x=(PQ-R)/3,\qquad
 xR_x=(PR-Q^2)/2;
\]
see \cite[Equations (1.1)--(1.4), p.~1]{ErbanVanGorder2021}.
Here $x=e^{2\pi i\tau}$, so
$d/d\tau=2\pi i x\,d/dx$.
For $E=\pi^2E_2/3$, $g_2=4\pi^4E_4/3$ and
$g_3=8\pi^6E_6/27$, this gives the exact normalization
\[
 \begin{split}
 E'&=\frac{i}{2\pi}(E^2-g_2/12),\\
 g_2'&=\frac{i}{\pi}(2Eg_2-3g_3),\qquad
 g_3'=\frac{i}{\pi}(3Eg_3-g_2^2/6).
 \end{split}
\]

\begin{proposition}[fourth order at the maximum]\label{ari:C.3.1}
\StatementRecord{ari:C.3.1}{proved}{Historical events: Claude Fable 5.2 (Anthropic), effort high, 16 November 2026: carried as background without fault found; used and rechecked in this run. GPT-6 Astra (OpenAI), ChatGPT Pro mode, effort Pro, S6\_09\_PAPERS, 17 November 2026: targeted proof check. Also checked in this writing. Pre-continuation writing events: GPT-6 Astra (OpenAI), ChatGPT Pro mode, effort Pro, S6\_10\_PAPERS, 18 November 2026: targeted proof check. Continuation mathematical and source checks: GPT-6 Astra (OpenAI), ChatGPT Pro mode, effort Pro, S6\_10\_PAPERS continuation, 19 November 2026: computation or formal certificate; independent re-proof; used and rechecked in this run. Continuation metadata events: GPT-6 Astra (OpenAI), ChatGPT Pro mode, effort Pro, S6\_10\_PAPERS continuation, 19 November 2026: used and rechecked in this run; status-vocabulary reconciliation; metadata reconciliation only; no new mathematical check. \emph{Continuation recorded proof scope.} Complete new mathematical derivation written before opening the target proof, followed by comparison of every new step. The fresh program \path{code/continuation_root/ari_C2_C3_fresh.py} passed 34 exact identities by a differential-polynomial and real-polynomial method. Ramanujan, theta and period inputs were opened with their exact normalizations; the renewed credit search and positive antecedents are recorded in the continuation comparison.}{the hypotheses, notation and definitions stated here; the explicit constructions and calculations in the proof}

On either root branch at $\tau=i$, let
$h=\mu-(1-\tau)/2$ and $\delta=\tau-i$. Then
\begin{align*}
 \tau&=i-\frac{3\pi i}{2}h^2-\frac{9\pi^2i}{8}h^4+O(h^6),\\
 h^2&=\frac{2i}{3\pi}\delta+\frac1{3\pi}\delta^2+O(\delta^3),\\
 \mu&=\frac{1-\tau}{2}
  \mathbin{\pm}\left(\frac{2i\delta}{3\pi}\right)^{1/2}
          \left(1-\frac{i}{4}\delta+O(\delta^2)\right).
\end{align*}
The square root is taken holomorphically on the chosen root chart.
Furthermore
\begin{align}\label{ari:eq:cs-Li}
 \log N={}&\tfrac12\log2-\frac\pi2|h|^2
   +\frac{3\pi^2}{8}|h|^4
   -\frac{3\pi^2}{4}|h|^2\operatorname{Re}(h^2)
   -\frac{5G}{192}\operatorname{Re}(h^4)+O(|h|^6),\\
 D_{\rm part}={}&D_* -\frac2\pi|\delta|
  +\frac1\pi\bigl(|\delta|^2+|\delta|\operatorname{Im}\delta\bigr)
  +\frac{5G}{72\pi^3}\operatorname{Re}(\delta^2)
  +O(|\delta|^3).\nonumber
\end{align}
In particular the leading terms agree with the boundary expansion in Lemma~\ref{push:src:D.2.1}.
\end{proposition}


\begin{proof}
\begin{enumerate}[label=\arabic*.,leftmargin=*]
\item Write $e(\tau)=\wp_\tau((1-\tau)/2)$ and
$E=\pi^2E_2/3$. The half-period Taylor series is
\[
 \wp_\tau((1-\tau)/2+h)
 =e+(3e^2-g_2/4)h^2+(3e^3-g_2e/4)h^4+O(h^6).
\]
The exact half-period coefficients used in Lemma~\ref{push:src:D.2.1} are
$e(i)=0$, $e'(i)=iG/(6\pi)$. This proves the claimed identity
$3e'(i)=(2i/\pi)G/4$. For clarity its closed form follows from
$e=\pi^2(\vartheta_2^4-\vartheta_4^4)/3$ and the heat equation:
if $A=\vartheta_2(i)^4=\vartheta_4(i)^4$, then
$G=4\pi^4A^2$ and $e'(i)=2\pi^3iA^2/3$.
For the classical closed value, \cite[Equation 23.5.3]{DLMFTheta2026}
gives $e_1=\Gamma(1/4)^4/(32\pi\omega_1^2)$,
$e_2=0$ and $e_3=-e_1$.  The lattice identification is
$2\omega_1=1$, $2\omega_3=i$; hence
$G=-4(e_1e_2+e_2e_3+e_3e_1)=4e_1^2
=\Gamma(1/4)^8/(16\pi^2)$.

\item Ramanujan's identities give $g_2'(i)=2iG$. Differentiating
$e(-1/\tau)=\tau^2 e(\tau)$ twice at $i$ gives
$e''(i)=3i e'(i)=-G/(2\pi)$. Inserting
$\delta=A_2h^2+A_4h^4+O(h^6)$ into the displayed half-period series
therefore gives
\[
 A_2=-3\pi i/2,\qquad
 A_4=\frac{g_2'A_2/4-e''A_2^2/2}{e'}=-9\pi^2i/8.
\]
The implicit-function theorem applies because $e'(i)\ne0$;
evenness in $h$ gives the stated remainder. Series inversion gives
the formulas for $h^2$ and $\mu$.

\item Let $H(h,\tau)=\log(\vartheta_3(h,\tau)/\eta(\tau))$. The exact
identity from the half-period translation is
$\log N=\operatorname{Re}H-\pi(\operatorname{Im}h)^2/
\operatorname{Im}\tau$. At $(0,i)$ its relevant Taylor data are
\[
 \begin{gathered}
 H_0=\tfrac12\log2,\quad H_\tau=0,\quad
 [h^2]H=-\pi/2,\quad [h^4]H=G/48,\\
 \partial_\tau[h^2]H=-\pi i/4-iG/(16\pi),\qquad
 H_{\tau\tau}=-G/(24\pi^2).
 \end{gathered}
\]
Indeed $\partial_h^2\log\vartheta_3=-\wp(h+(1-\tau)/2)-E$;
the heat equation at $h=0$ gives $H_\tau=ie/(4\pi)$, and
$E'(i)=\pi i/2-iG/(24\pi)$. These identities prove every entry.
After substitution $\tau=\tau(h)$, the pure $h^4$ coefficient of
$H$ is $-3\pi^2/8-5G/192$. Expanding
$\operatorname{Im}\tau=1-(3\pi/2)\operatorname{Re}(h^2)+O(|h|^4)$
in the Gaussian gives \eqref{ari:eq:cs-Li}. Finally
\[
 |h|^2=\frac{2}{3\pi}|\delta|
          \left(1+\tfrac12\operatorname{Im}\delta+O(|\delta|^2)\right)
\]
and substitution into the quartic expression, followed by
multiplication by $6/\pi$, gives the formula for $D_{\rm part}$.
The fresh exact differential-polynomial calculation in
\path{code/continuation_root/ari_C2_C3_fresh.py}
checks all these coefficients and the real quartic polynomial.
Its exact output is retained beside the program; the earlier theta
evaluation diagnostics remain in the computation record.
\end{enumerate}
\end{proof}

\begin{proposition}[the equianharmonic point is neither an extremum]\label{ari:C.3.2}
\StatementRecord{ari:C.3.2}{proved}{Historical events: Claude Fable 5.2 (Anthropic), effort high, 15 November 2026: carried as background without fault found. GPT-6 Astra (OpenAI), ChatGPT Pro mode, effort Pro, S6\_09\_PAPERS, 17 November 2026: targeted proof check. Also checked in this writing. Pre-continuation writing events: GPT-6 Astra (OpenAI), ChatGPT Pro mode, effort Pro, S6\_10\_PAPERS, 18 November 2026: targeted proof check. Continuation mathematical and source checks: GPT-6 Astra (OpenAI), ChatGPT Pro mode, effort Pro, S6\_10\_PAPERS continuation, 19 November 2026: computation or formal certificate; independent re-proof; used and rechecked in this run. Continuation metadata events: GPT-6 Astra (OpenAI), ChatGPT Pro mode, effort Pro, S6\_10\_PAPERS continuation, 19 November 2026: used and rechecked in this run; status-vocabulary reconciliation; metadata reconciliation only; no new mathematical check. \emph{Continuation recorded proof scope.} Complete new mathematical derivation written before opening the target proof, followed by comparison of every new step. The fresh program \path{code/continuation_root/ari_C2_C3_fresh.py} passed 34 exact identities by a differential-polynomial and real-polynomial method. Ramanujan, theta and period inputs were opened with their exact normalizations; the renewed credit search and positive antecedents are recorded in the continuation comparison.}{the hypotheses, notation and definitions stated here; the explicit constructions and calculations in the proof}

Set $T_0=\sqrt3/2$, $\mu_0=(2-\rho)/3$, and
$p=\wp'_\rho(\mu_0)$. Then $p\ne0$, $p^2=-g_3(\rho)$, and, for
$\delta=\tau-\rho$ on the order-three root chart,
\begin{align*}
 \mu&=\frac{2-\tau}{3}
      +\frac{3p}{16\pi^2}\delta^2
      +\frac{3ip}{32\pi^2T_0}\delta^3+O(\delta^4),\\
 t&=\frac{ip^2}{\pi^3}\delta^3
        \left(1+\frac{3i}{2T_0}\delta+O(\delta^2)\right),\\
 \log N&=\tfrac14\log3+\frac7{192}\operatorname{Re}t
       -\frac{9\pi}{512T_0|p|^{2/3}}|t|^{4/3}
       +O(|t|^{5/3}).
\end{align*}
Thus
\[
 D_{\rm part}(t)=-\frac3{2\pi}\log48
       +\frac7{32\pi}\operatorname{Re}t
       -\frac{27}{256T_0|p|^{2/3}}|t|^{4/3}
       +O(|t|^{5/3}).
\]
The point $t=0$ is neither a local maximum nor a local minimum.
\end{proposition}


\begin{proof}
\begin{enumerate}[label=\arabic*.,leftmargin=*]
\item We rederive the differentiation in Proposition~\ref{push:src:D.4.4}: put
$\ell=\partial_z\log\vartheta_1(\mu,\tau)$,
$a=\operatorname{Im}\mu/T$, $Z=\ell+2\pi ia$, and
$R=\mu'-a$. The theta heat equation and
$\wp=-\partial_z\ell-E$ give on the zero locus
\[
 \mu'=\frac{i\ell}{2\pi}-\frac{ig_2}{6\pi\wp'},\qquad
 R=\frac{iZ}{2\pi}-\frac{ig_2}{6\pi\wp'},\qquad
 \partial_\tau\log N=\frac{Z}{8\pi i}(Z+4\pi iR).
\]
These equations give the factored formula of Proposition~\ref{push:src:D.4.4}
with its stated normalizations. At $\rho$ the unit endomorphism fixing the
nonzero order-three point forces the gradient of $U_\rho$ there to
vanish, so $Z=0$. Explicitly the unit $\rho-1$ fixes $\mu_0$ modulo
the lattice and has no fixed nonzero real covector. Hence
$a=\mu'=-1/3$ and $\ell=2\pi i/3$.

\item Along the zero locus, a further heat-equation differentiation gives
\[
 \ell'=\frac{i\wp'}{4\pi}+\frac{iEg_2}{6\pi\wp'}.
\]
The Ramanujan derivatives at $\rho$ are
\[
 g_2'=\frac{3ip^2}{\pi},\quad
 g_2''=\frac{5i}{T_0}g_2',\quad
 \frac{(\wp')'}{\wp'}=\frac{3i}{2T_0},\quad E=\frac\pi{T_0}.
\]
They give $\mu''=3p/(8\pi^2)$,
$\mu'''=9ip/(16\pi^2T_0)$ and hence the asserted $\mu$ expansion.
Expansion of $t=g_2^3/(g_2^3-27g_3^2)$ gives the displayed $t$
formula. Since $\partial_\tau a=R/(2iT)$, one has
\[
 Z=\frac{ip}{4\pi}\delta+O(|\delta|^2),\qquad
 R=\frac{3p}{8\pi^2}\delta+O(|\delta|^2).
\]
The derivative formula now gives the harmonic cubic
$\operatorname{Re}(7ip^2\delta^3/(192\pi^3))=(7/192)
\operatorname{Re}t+O(|\delta|^4)$.

\item For the next term, direct differentiation of the Gaussian gives
\[
 \partial_\tau\partial_{\bar\tau}\log N
       =-\frac\pi{2T}|R|^2
       =-\frac{9|p|^2}{128\pi^3T_0}|\delta|^2+O(|\delta|^3).
\]
After the harmonic cubic has been written as a function of $t$, the
order-three stabilizer permits only a radial term in degree four:
the actual stabilizer map is $\gamma\tau=1-1/\tau$ and the zero
branch transforms as
$\mu(\gamma\tau)=(1-\mu(\tau))/\tau$.
The theta and eta transformations cited above cancel against the
Gaussian factor, so $N(\gamma\tau)=N(\tau)$.
Thus
in the linearizing coordinate
$(\tau-\rho)/(\tau-\bar\rho)$ invariant monomials have exponents
whose difference is divisible by three. The coefficient of
$|\delta|^4$ is therefore $-9|p|^2/(512\pi^3T_0)$, since
$\partial_\delta\partial_{\bar\delta}|\delta|^4=4|\delta|^2$.
Converting $|\delta|^4$ to $|t|^{4/3}$ proves the expansion, including
the remainder. Its nonzero linear term in $\operatorname{Re}t$
has both signs, which proves the last assertion. In particular
superharmonicity at this root point does not make it a local maximum.
\end{enumerate}
\end{proof}



\subsection{CM and torsion}
\begin{cited}[no smooth CM-torsion parameter]\label{ari:C.5.1}
\StatementRecord{ari:C.5.1}{proved}{Historical events: Claude Opus 5.5 (Anthropic), effort max, sweep 1, scope s6, 18 November 2026: used and rechecked in this run. GPT-6 Astra (OpenAI), ChatGPT Pro mode, effort Pro, S6\_09\_PAPERS, 17 November 2026: targeted proof check. This writing records metadata reconciliation only; no new mathematical check. Continuation metadata events: GPT-6 Astra (OpenAI), ChatGPT Pro mode, effort Pro, S6\_10\_PAPERS continuation, 19 November 2026: used and rechecked in this run; status-vocabulary reconciliation; metadata reconciliation only; no new mathematical check. \emph{Pre-continuation recorded proof scope.} Current-writing verification remains in progress. The earlier checks named above retain their recorded scope.}{the hypotheses, notation and definitions stated here; \citet[Theorem~4.15]{BourdonClarkStankewicz2017}; \citet[Theorem~20.12 and Corollary~20.13]{Sutherland2023CM}; \cref{app:rederive:ari-C.5.1}}

There is no $t\in\mathbb C\setminus\{0,1\}$ for which $E_t$ has
complex multiplication and $P_t$ is torsion. On the smooth models
at $t=0,1$ after base change of orders $3,4$, respectively, the
limiting points have orders $3,2$. Thus these are the only CM
torsion points on the completed family.
\end{cited}
The expanded source derivation is retained in the proof compendium, Paper~I, label \nolinkurl{ari:C.5.1}.




This applies \citet[Theorem~4.15]{BourdonClarkStankewicz2017} and \citet[Theorem~20.12 and Corollary~20.13]{Sutherland2023CM}; the exact finite enumeration is retained. A complete re-derivation is given in Appendix~\ref{app:rederive:ari-C.5.1}.



\section{Global heights and the eleven rational CM fibres}

We use the equation
\[
 E_t:Y^2=X^3-\frac34t^3(t-1)X+\frac14t^4(t-1)^2,
 \qquad P_t=(0,t^2(t-1)/2).
\]
Its discriminant is $27t^8(t-1)^3$. Our canonical height is
$\widehat h_O(P)=\frac12\lim_{n\to\infty}n^{-2}h(x(nP))$.
In particular the bilinear Mordell--Weil self-pairing used in the
supplied surface calculation is $2\widehat h_O$.
For a finite place $v$ above the rational prime $p$, we normalize
$|p|_v=p^{-1}$ and set $n_v=[K_v:\Q_p]$. At an archimedean
place $n_v$ is one or two, according as the place is real or complex.

\begin{proposition}[the global formula]\label{ari:C.6.1}
\StatementRecord{ari:C.6.1}{proved}{Historical events: GPT-6 Astra (OpenAI), ChatGPT Pro mode, effort Pro, S6\_09\_PAPERS, 17 November 2026: targeted proof check. Also checked in this writing. Pre-continuation writing events: GPT-6 Astra (OpenAI), ChatGPT Pro mode, effort Pro, S6\_10\_PAPERS, 18 November 2026: targeted proof check; used and rechecked in this run. Continuation metadata events: GPT-6 Astra (OpenAI), ChatGPT Pro mode, effort Pro, S6\_10\_PAPERS continuation, 19 November 2026: used and rechecked in this run; status-vocabulary reconciliation; metadata reconciliation only; no new mathematical check. \emph{Pre-continuation recorded proof scope.} Read complete printed arithmetic proofs and re-derived the exact scaling/local-height equations and the genus-trace square denominator argument. Rank, saturation, directed L-derivative replays and precise external theorem matches remain explicit gates; no proof verdict is awarded to those computational claims by this read. Opened DeMarco–Mavraki1701.07947v2 Section2.1 pp5–6 and matched the canonical/local height normalization, continuity/pole properties and duplication discriminant term1/4 log|Delta|. This closes this source-normalization gate.}{the hypotheses, notation and definitions stated here; \cref{ari:C.7.1}; \citet[Section 2.1]{demarco2017}}

For every number field $K$, every $t\in K\setminus\{0,1\}$,
and the section just specified,
\begin{equation}\label{ari:eq:ch-global}
 \sum_{\sigma:K\hookrightarrow\C}\log N(P_{\sigma(t)})
 =\sum_{v\nmid\infty}n_v\lambda_v(P_t)
          -[K:\Q]\widehat h_O(P_t).
\end{equation}
Here the local heights are invariant under change of Weierstrass
model and satisfy, when $nP\ne O$,
\begin{equation}\label{ari:eq:ch-local-mult}
 \lambda_v(nP)=n^2\lambda_v(P)-\log|\psi_n(P)|_v
                   +\frac{n^2-1}{12}\log|\Delta|_v.
\end{equation}
At a complex embedding $\lambda_v(P)=-\log N(P)$; complex
embeddings are counted separately in \eqref{ari:eq:ch-global}.
For rational $t$, put $A_t=\exp(\sum_p\lambda_p(P_t))$.
Then $A_t$ is a positive real algebraic number and
\[
 N(P_t)=A_t e^{-\widehat h_O(P_t)},\qquad
 D_{\rm part}(t)=\frac6\pi
       \left(\log\frac{A_t}{2\sqrt3}-\widehat h_O(P_t)\right).
\]
Thus $N(P_t)$ is algebraic if and only if
$e^{\widehat h_O(P_t)}$ is algebraic.
\end{proposition}


\begin{proof}
\begin{enumerate}[label=\arabic*.,leftmargin=*]
\item The sigma expression for the archimedean local height is
\[
 \lambda_v(P)=-\log|\sigma(u)|
       +\frac12\Re(u\eta(u))-\frac1{12}\log|\Delta|.
\]
The sigma--theta relation and the Legendre relation identify this
with $-\log N$ exactly, including its constant; this is also the
calculation in Statement \ref{ari:C.7.1}. The division identity
$\sigma(nu)=(-1)^{n-1}\psi_n(P)\sigma(u)^{n^2}$ gives
\eqref{ari:eq:ch-local-mult}. At finite places the same identity is
the defining multiplication law of the canonical N\'eron function.
Its normalization is uniquely determined by continuity away from
$O$, its logarithmic pole at $O$, and the multiplication law.
These are the three properties printed in
\citet[Section 2.1]{demarco2017}; their duplication law includes
the discriminant term $\frac14\log|\Delta|_v$.
Under a coordinate scaling the changes in the sigma or formal
parameter term and in $\frac1{12}\log|\Delta|_v$ cancel.

\item Sum the duplication law over places. The product formula cancels
the division polynomial and discriminant terms, so that the sum
$H(P)=[K:\Q]^{-1}\sum_vn_v\lambda_v(P)$ satisfies
$H(2P)=4H(P)$. The pole normalization implies
$H(P)=\frac12h(x(P))+O_E(1)$. Iterating multiplication by two
therefore identifies $H$ with $\widehat h_O$. Moving the
archimedean terms to the other side gives \eqref{ari:eq:ch-global}.
For a finite place, a multiple of $P$ lies in the identity component
of a minimal model. There the local height is a rational multiple
of $\log p$; the multiplication law shows the same for $P$.
Only finitely many places contribute. This proves the assertion
about $A_t$ and both of its consequences.
\end{enumerate}
\end{proof}

\begin{proposition}[exact finite local-height calculation]\label{ari:C.6.2}
\StatementRecord{ari:C.6.2}{proved}{Historical events: GPT-6 Astra (OpenAI), ChatGPT Pro mode, effort Pro, S6\_09\_PAPERS, 17 November 2026: targeted proof check. Also checked in this writing. Pre-continuation writing events: GPT-6 Astra (OpenAI), ChatGPT Pro mode, effort Pro, S6\_10\_PAPERS, 18 November 2026: targeted proof check. Continuation metadata events: GPT-6 Astra (OpenAI), ChatGPT Pro mode, effort Pro, S6\_10\_PAPERS continuation, 19 November 2026: used and rechecked in this run; status-vocabulary reconciliation; metadata reconciliation only; no new mathematical check. \emph{Pre-continuation recorded proof scope.} Read complete printed arithmetic proofs and re-derived the exact scaling/local-height equations and the genus-trace square denominator argument. Rank, saturation, directed L-derivative replays and precise external theorem matches remain explicit gates; no proof verdict is awarded to those computational claims by this read.}{the hypotheses, notation and definitions stated here; \cref{ari:eq:ch-local-mult}}

Let $E/\Q$ be a global minimal Weierstrass model. If $mP$ is
in the identity component at $p$, and $mP\ne O$, then
\begin{equation}\label{ari:eq:ch-local-exact}
 \frac{\lambda_p(P)}{\log p}
 =\frac{v_p(\Delta)}{12}
   +\frac{\frac12\max(0,-v_p(x(mP)))-v_p(\psi_m(P))}{m^2}.
\end{equation}
For $m=1$ we read $\psi_1=1$. This formula, with exact rational
coordinates and valuations, gives all the finite contributions
in the tables below.
\end{proposition}


\begin{proof}
\begin{enumerate}[label=\arabic*.,leftmargin=*]
\item On the identity component of a minimal model the canonical
N\'eron function is
$\frac12\log\max(1,|x|_p)-\frac1{12}\log|\Delta|_p$.
This follows either from the formal parameter at $O$ followed
by the group law on the smooth locus, or from the uniqueness
properties used above. Substitute it for $\lambda_p(mP)$ in
\eqref{ari:eq:ch-local-mult} and divide by $m^2\log p$.
The contribution $v_p(\Delta)/(12m^2)$ combines with the
discriminant term to give $v_p(\Delta)/12$.
Membership in the identity component is checked by reduction
to the nonsingular locus of the minimal cubic. Thus the test is
exact also at $2$ and $3$; it does not use a formula restricted
to residue characteristic at least five.
\end{enumerate}
\end{proof}

\begin{proposition}[the CM models and their height identities]\label{ari:C.6.3}
\StatementRecord{ari:C.6.3}{proved}{Historical events: GPT-6 Astra (OpenAI), ChatGPT Pro mode, effort Pro, S6\_09\_PAPERS, 17 November 2026: targeted proof check. Also checked in this writing. Pre-continuation writing events: GPT-6 Astra (OpenAI), ChatGPT Pro mode, effort Pro, S6\_10\_PAPERS, 18 November 2026: computation or formal certificate; targeted proof check. Continuation metadata events: GPT-6 Astra (OpenAI), ChatGPT Pro mode, effort Pro, S6\_10\_PAPERS continuation, 19 November 2026: used and rechecked in this run; status-vocabulary reconciliation; metadata reconciliation only; no new mathematical check. \emph{Pre-continuation recorded proof scope.} Read complete printed arithmetic proofs and re-derived the exact scaling/local-height equations and the genus-trace square denominator argument. Rank, saturation, directed L-derivative replays and precise external theorem matches have their separately recorded scopes. Reran C\_cm\_height\_arb.py at 1152 bits with 350 Fourier terms. All eleven CM j intervals certify the printed integers; all prescribed uniformizing root brackets are proved by opposite interval signs, and all theta-height radii are below 1e-236. Reran the exact model/local-height programs in PARI2.17.2 and the unchanged eclib20250122 saturation code. All eleven displayed models and local coefficients match; ranks are exactly2 for rows1,8 and exactly1 for the remaining nine. All nine saturation indices are1, with no remaining primes. The full actual outputs and source hashes are retained. Opened PARI ellrank/ellheight normalization documentation and Cremona--Siksek Theorem1, author-PDFp8, for the algorithmic scope.}{the hypotheses, notation and definitions stated here; \cref{ari:eq:ch-local-exact}}

The eleven rational CM fibres have the global minimal models,
section coordinates, component classes, and finite local heights
in the following tables. On each model $a_1=a_2=0$, and we
write $[a_3,a_4,a_6]$ for
$y^2+a_3y=x^3+a_4x+a_6$. The displayed point is $(0,y_P)$.
The sign of $P$, which is immaterial to all these heights,
is chosen by the recorded change of variables.
The constants $A_t$ in increasing row order are
\[
 \sqrt6,\quad 2^{2/3}3^{1/4},\quad 2^{3/4},\quad
 2,\quad 2,\quad \sqrt3,\quad 1,\quad 3^{5/12},\quad
 1,\quad1,\quad1.
\]
\end{proposition}


\begin{proof}
\begin{enumerate}[label=\arabic*.,leftmargin=*]
\item We substitute the eleven exact values $t=j/1728$, and perform
the integral minimal-model transformations and Tate reduction
algorithm. The transformations and resulting invariants are
printed in \path{code/C_cm_models.out}; the exact valuation
computations are in \path{code/C_cm_local_heights.out}.
For each bad prime the program computes $mP$ by the rational
group law, tests the two partial derivatives of the minimal
cubic, and evaluates $\psi_m$ by its exact recursion. Formula
\eqref{ari:eq:ch-local-exact} then proves each entry. Outside the
listed primes the discriminant is a unit and $P$ is integral,
so the local height is zero. Multiplying $p^{\lambda_p/\log p}$
gives the listed constants $A_t$.

\item We specify the component labels without relying on an
unrecorded ordering of a resolution graph. The label $0$
means the identity component, and $c_P$ is the class represented
by the explicitly displayed section $P$. At an odd prime of
type $I_0^*$ this nonidentity class is the outer component
indexed by the root zero of the residual Tate cubic; the
scaling $x=pX$ identifies that root. When there are two
rational components, $c_P$ is their unique nonzero class.
The only remaining ambiguity is row 5 at $2$, where the
component group is $(\Z/2)^2$. Here let $T=(-14,0)$;
we label it with the basis $(c_P,c_T)$. Indeed
$P+T=(18,64)$ after the corresponding choice of sign,
and $P,T,P+T$ all reduce to the singular point, whereas
$2P=(25/4,99/8)$ reduces to the identity component.
Thus the three nonzero classes are explicitly separated.

\item The independent analytic check uses the exact quadratic CM
period ratio, solves $\wp(\mu)=0$ by its convergent Fourier
series, and evaluates the theta product at 270-digit working
precision. The value of $j$ agrees with its prescribed integer
to absolute error less than $10^{-245}$. For every row the
direct value of $\log N$ agrees with
$\log A_t-\widehat h_O(P_t)$ to more than 245 decimal places.
These numerical comparisons test the exact normalization;
the algebraic local-height entries were already proved by
integer and rational arithmetic. The full precision values
are in \path{records/C_cm_theta_pslq.json}.
\end{enumerate}
\end{proof}

\begingroup\footnotesize\setlength{\tabcolsep}{4pt}
\begin{longtable}{r r r r r r}
\toprule Row & $j$ & $a_3$ & $a_4$ & $a_6$ & $y_P$\\\midrule\endhead
1 & 8000 & 0 & -5880 & 153664 & 392\\
2 & 54000 & 0 & -7260 & 234256 & 484\\
3 & 287496 & 0 & -2156 & 38416 & 196\\
4 & 16581375 & 0 & -7732620 & 8275904784 & 90972\\
5 & -3375 & 0 & -140 & 784 & -28\\
6 & -32768 & 1 & -3234 & 72630 & -270\\
7 & -884736 & 1 & -38 & 90 & -10\\
8 & -12288000 & 1 & -1920270 & 1024288020 & -32005\\
9 & -884736000 & 1 & -379260 & 89898842 & -9482\\
10 & -147197952000 & 1 & -347045930 & 2488447882842 & -1577482\\
11 & -262537412640768000 & 1 & -75065352465446420 & 7916029533155605844266842 & -2813543945482\\
\bottomrule\end{longtable}
\begin{longtable}{r r l c c l}
\toprule Row & $p$ & Type & Component & Order & $\lambda_p/\log p$\\\midrule\endhead
1 & 2 & $III^*$ & $c_P$ & 2 & $\frac{1}{2}$\\
1 & 3 & $I_0^*$ & $0$ & 1 & $\frac{1}{2}$\\
1 & 7 & $I_0^*$ & $c_P$ & 2 & $0$\\
2 & 2 & $I_{4}^*$ & $c_P$ & 2 & $\frac{2}{3}$\\
2 & 3 & $III$ & $0$ & 1 & $\frac{1}{4}$\\
2 & 11 & $I_0^*$ & $c_P$ & 2 & $0$\\
3 & 2 & $I_{5}^*$ & $c_P$ & 2 & $\frac{3}{4}$\\
3 & 7 & $I_0^*$ & $c_P$ & 2 & $0$\\
4 & 2 & $I_{8}^*$ & $c_P$ & 2 & $1$\\
4 & 3 & $I_0^*$ & $c_P$ & 2 & $0$\\
4 & 7 & $III$ & $c_P$ & 2 & $0$\\
4 & 19 & $I_0^*$ & $c_P$ & 2 & $0$\\
5 & 2 & $I_{8}^*$ & $c_P$ & 2 & $1$\\
5 & 7 & $III$ & $c_P$ & 2 & $0$\\
6 & 3 & $I_0^*$ & $0$ & 1 & $\frac{1}{2}$\\
6 & 7 & $I_0^*$ & $c_P$ & 2 & $0$\\
6 & 11 & $III$ & $c_P$ & 2 & $0$\\
7 & 19 & $III$ & $c_P$ & 2 & $0$\\
8 & 3 & $IV$ & $0$ & 1 & $\frac{5}{12}$\\
8 & 11 & $I_0^*$ & $c_P$ & 2 & $0$\\
8 & 23 & $I_0^*$ & $c_P$ & 2 & $0$\\
9 & 3 & $I_0^*$ & $c_P$ & 2 & $0$\\
9 & 7 & $I_0^*$ & $c_P$ & 2 & $0$\\
9 & 43 & $III$ & $c_P$ & 2 & $0$\\
10 & 7 & $I_0^*$ & $c_P$ & 2 & $0$\\
10 & 31 & $I_0^*$ & $c_P$ & 2 & $0$\\
10 & 67 & $III$ & $c_P$ & 2 & $0$\\
11 & 7 & $I_0^*$ & $c_P$ & 2 & $0$\\
11 & 11 & $I_0^*$ & $c_P$ & 2 & $0$\\
11 & 19 & $I_0^*$ & $c_P$ & 2 & $0$\\
11 & 127 & $I_0^*$ & $c_P$ & 2 & $0$\\
11 & 163 & $III$ & $c_P$ & 2 & $0$\\
\bottomrule\end{longtable}
\endgroup
\noindent The next table prints 52 significant digits. The last column is both the directly computed $\log N$ and the independently computed $\log A_t-\widehat h_O(P_t)$; the agreement is more precise than the digits printed.
\begingroup\small
\begin{longtable}{r l}
\toprule Row & Height and archimedean logarithm\\\midrule\endhead
1 & $\widehat h_O=0.68444758838339190721754413638128047178526153154571$\\
 & $\log N=0.2114321462306355931886945428090706645762338145457924$\\[3pt]
2 & $\widehat h_O=0.6417998032172362080432024005106271224631326545214062$\\
 & $\log N=0.09495138932308808775043032302545534908240674117445132$\\[3pt]
3 & $\widehat h_O=0.5335302122611885050129127004864161587766141055417204$\\
 & $\log N=-0.01366982684122952294998860939278373271998900477152897$\\[3pt]
4 & $\widehat h_O=0.9940668685687274300038154462317330832788241518412543$\\
 & $\log N=-0.300919688008782120586583324773556515203324017480999$\\[3pt]
5 & $\widehat h_O=0.4673656771814712800075202012628467836356203437033082$\\
 & $\log N=0.2257815033784740294097119201953297844398797906569471$\\[3pt]
6 & $\widehat h_O=0.4286091723813209970260194687110167110834518087280487$\\
 & $\log N=0.120696971952733848671603149750246141240293470183326$\\[3pt]
7 & $\widehat h_O=0.09064583701789963707894077069188164993051948122321821$\\
 & $\log N=-0.09064583701789963707894077069188164993051948122321821$\\[3pt]
8 & $\widehat h_O=0.7366208256267509672496883390906560396314066510934334$\\
 & $\log N=-0.2788657053483719291683361570396036626949522520006211$\\[3pt]
9 & $\widehat h_O=0.6018296198896332298783588845164397135769872973598109$\\
 & $\log N=-0.6018296198896332298783588845164397135769872973598109$\\[3pt]
10 & $\widehat h_O=1.002649286144114870320961646648539589847474352238484$\\
 & $\log N=-1.002649286144114870320961646648539589847474352238484$\\[3pt]
11 & $\widehat h_O=2.165489246796061077105970244492001901750597685962121$\\
 & $\log N=-2.165489246796061077105970244492001901750597685962121$\\[3pt]
\bottomrule\end{longtable}\endgroup


\begin{proposition}[the maximum as a height inequality]\label{ari:C.6.4}
\StatementRecord{ari:C.6.4}{proved}{Historical events: Claude Fable 5.2 (Anthropic), effort high, 15 November 2026: carried as background without fault found. GPT-6 Astra (OpenAI), ChatGPT Pro mode, effort Pro, S6\_09\_PAPERS, 17 November 2026: targeted proof check. Also checked in this writing. Pre-continuation writing events: GPT-6 Astra (OpenAI), ChatGPT Pro mode, effort Pro, S6\_10\_PAPERS, 18 November 2026: targeted proof check; used and rechecked in this run. Continuation metadata events: GPT-6 Astra (OpenAI), ChatGPT Pro mode, effort Pro, S6\_10\_PAPERS continuation, 19 November 2026: used and rechecked in this run; status-vocabulary reconciliation; metadata reconciliation only; no new mathematical check. \emph{Pre-continuation recorded proof scope.} Read complete printed arithmetic proofs and re-derived the exact scaling/local-height equations and the genus-trace square denominator argument. Rank, saturation, directed L-derivative replays and precise external theorem matches remain explicit gates; no proof verdict is awarded to those computational claims by this read. Opened DeMarco–Mavraki Theorem1.1 p2, Theorems2.2–2.3 pp7–8 and the metric/proof loci pp15–16,23; matched the positive generic height1/12 and the finite-place/continuous-remainder application. Other explicitly recorded CM rank or directed-numerical gates are unchanged. Opened Silverman author slides21–22 (PDF26–27), matching the uniform height bound and specialization limit with their hypotheses. The exact DeMarco–Mavraki application and its bounded local remainder remain the operative proof. The original1983Silverman and Tate articles remain explicitly unverified background records.}{the hypotheses, notation and definitions stated here; \cref{ari:C.8.5}; \cref{ari:eq:ch-global}; \citet{demarco2017}; \citet[pp.~183--184]{callSilverman1993}; \citet{silverman2006heights}}

For $t\in K\setminus\{0,1\}$,
\begin{equation}\label{ari:eq:ch-bound}
 \widehat h_O(P_t)\geq
 \frac1{[K:\Q]}\sum_{v\nmid\infty}n_v\lambda_v(P_t)
                    -\frac12\log2.
\end{equation}
The inequality is strict on this smooth base. Over $\Q$ it
is $\widehat h_O(P_t)\geq\log A_t-\frac12\log2$.
The generic height is $\widehat h_O(P)=1/12$, and the
variation divisor is $D_P=(\infty)/12$. Consequently
\[
 \widehat h_O(P_t)=\frac1{12}h(t)+O(1)
 \qquad(t\in\overline\Q\setminus\{0,1\}),
\]
where the constant is uniform over algebraic $t$ and is for
the fixed family and section.
\end{proposition}


\begin{proof}
\begin{enumerate}[label=\arabic*.,leftmargin=*]
\item Apply the analytic maximum theorem \ref{ari:C.8.5} at every embedding
and substitute into \eqref{ari:eq:ch-global}. Equality would
require $j(\sigma(t))=1728$ for every embedding, hence $t=1$,
which has been excluded.

\item For the generic local heights at $0$ and $1$, the minimal
discriminant orders are $8$ and $3$, and the section meets
the nonidentity component of types $IV^*$ and $III$.
Their local corrections are $2/3$ and $1/4$, respectively,
so their normalized local heights are zero. At infinity,
put $s=1/t$, $X=s^{-2}x$, $Y=s^{-3}y$. The new equation is
\[
 y^2=x^3-\tfrac34(1-s)x+\tfrac14(1-s)^2,
 \qquad P=(0,(1-s)/2),\qquad\Delta=27s(1-s)^3.
\]
At $s=0$ the node is $(1/2,0)$, whereas $P=(0,1/2)$ is
nonsingular. Its height contribution is $1/12$.
At all other function-field places the point is integral
and the discriminant is a unit. This proves
$D_P=(\infty)/12$ and the generic height.

\item Theorem 1.1 of \citet{demarco2017} applies with $K=\Q$,
$B=\mathbb P^1$, the displayed elliptic surface and rational
section, since the generic height is positive. It equips
this divisor with continuous adelic semipositive metrics
whose associated height is exactly $\widehat h_O(P_t)$.
Comparing with the usual height for $(\infty)/12$ gives the
bounded difference. This is the specialization here of the
variation theorem of Tate, with the local refinements of
Silverman reproduced as Theorems 2.2 and 2.3 in that paper:
outside finitely many places the local correction is zero,
and at each remaining place it extends continuously after
subtracting its prescribed logarithmic singularity.
Our inequality is an explicit archimedean estimate with
constant $\frac12\log2$; the general bounded-variation
statement alone does not supply that constant.

\item We spell out the comparison with the opened primary author sources.
Theorem 4.1 of \citet[pp.~183--184]{callSilverman1993} concerns a
smooth projective base curve $T$, a degree-one base height $h_T$,
a section, and a fibrewise morphism whose pullback multiplies the
chosen divisor class by a real number $\alpha>1$. It gives
\[
 \lim_{h_T(t)\to\infty}
 \frac{\widehat h_t(P_t)}{h_T(t)}
       =\widehat h_{\Q(T)}(P).
\]
Here take $T=\mathbb P^1$, the class of $(O)$, and the generic
morphism $[2]$, with $[2]^*(O)\sim4(O)$ and $\alpha=4$.
The smooth fibres are precisely the complement of $0,1,\infty$;
the section extends on the regular projective elliptic surface.
All the hypotheses hold, and the limit is $1/12$.
The unnumbered remark immediately following that theorem states
Tate's stronger elliptic-family formula
$\widehat h_t(P_t)=h_{T,D_P}(t)+O_P(1)$, with
$\deg D_P=\widehat h_{\Q(T)}(P)$. This is an author-verified
reproduction of Tate's result, not a claim to have read the
inaccessible original scan \citep{tate1983variation}.
The same distinction applies to \citet{silverman1983specialization}:
the limit theorem is reproduced by the author on slide~22 of
\citet{silverman2006heights}, while slide~21 states the uniform
canonical-versus-Weil-height bound in terms of the base height.
For our explicit bounded remainder and local normalization, the
fully opened Theorems~1.1, 2.2 and 2.3 of
\citet{demarco2017}, applied above, remain the cited proof.
\end{enumerate}
\end{proof}

\begin{proposition}[rigorous high-precision height enclosures]\label{ari:C.6.5}
\StatementRecord{ari:C.6.5}{proved}{Historical events: GPT-6 Astra (OpenAI), ChatGPT Pro mode, effort Pro, S6\_09\_PAPERS, 17 November 2026: targeted proof check. Also checked in this writing. Pre-continuation writing events: GPT-6 Astra (OpenAI), ChatGPT Pro mode, effort Pro, S6\_10\_PAPERS, 18 November 2026: targeted proof check. Continuation metadata events: GPT-6 Astra (OpenAI), ChatGPT Pro mode, effort Pro, S6\_10\_PAPERS continuation, 19 November 2026: used and rechecked in this run; status-vocabulary reconciliation; metadata reconciliation only; no new mathematical check. \emph{Pre-continuation recorded proof scope.} Read complete printed arithmetic proofs and re-derived the exact scaling/local-height equations and the genus-trace square denominator argument. Rank, saturation, directed L-derivative replays and precise external theorem matches remain explicit gates; no proof verdict is awarded to those computational claims by this read.}{the hypotheses, notation and definitions stated here; \cref{ari:C.6.3}; \cref{ari:C.5.1}}

For every one of the eleven rational CM parameters in \ref{ari:C.6.3},
the files \path{records/C_cm_height_arb.json} and
\path{code/C_cm_height_arb.out} give certified enclosures of
$\log N(P)$ and $\widehat h_O(P)$ of radius less than
$10^{-236}$. These enclosures use no floating-point height
value as a proof step.
\end{proposition}


\begin{proof}
\begin{enumerate}[label=\arabic*.,leftmargin=*]
\item The CM discriminants, in the order of \ref{ari:C.6.3}, are
\[
 -8,-12,-16,-28,-7,-11,-19,-27,-43,-67,-163.
\]
Enumeration of primitive reduced forms
$(a,b,c)$, $|b|\le a\le c$, $b^2-4ac=D$, with the usual
nonnegative boundary convention, gives exactly one form of
each discriminant, always with $a=1$.
The range is finite and proved by $a\le\sqrt{|D|/3}$.
Thus the CM theorem used in \ref{ari:C.5.1} makes the corresponding
singular modulus an integer. The Fourier enclosures below
place it within $1/4$ of the specified $j$, proving its
identity, rather than only numerical closeness. Exact class
polynomials were also computed independently in
\path{code/C_cm_class_polynomials.gp} and its output.

\item Write $\tau=iT$ for the first four entries and
$\tau=1/2+iT$ for the other seven, with
\[
 (T_1,T_2,T_3,T_4)=(\sqrt2,\sqrt3,2,\sqrt7),\qquad
 T_k=\sqrt{|D_k|}/2\quad(k\ge5).
\]
Put $Q=e^{-2\pi T}$, $q=Q$ or $-Q$ accordingly,
$M=350$, and $Y=2/5$. On $0\le y\le Y$,
\begin{equation}\label{ari:eq:ch-wp-bracket}
 \frac{\wp(1/2-iy)}{\pi^2}
 =\frac1{\cosh^2\pi y}-\frac13+
 8\sum_{n\ge1}\frac{nq^n}{1-q^n}
                 \{1-(-1)^n\cosh(2\pi ny)\}.
\end{equation}
For $a=Qe^{2\pi Y}$, the absolute tail after $M$ in
$\wp$ is at most
\[
 \frac{16\pi^2}{1-Q}
 \frac{a^{M+1}\{M+1-Ma\}}{(1-a)^2}<10^{-875}.
\]
The proof is $1+\cosh(2\pi ny)\le2e^{2\pi nY}$ and
$\sum_{n>M}na^n=a^{M+1}(M+1-Ma)/(1-a)^2$.
Each previously computed decimal supplies only a rational
candidate centre $y_0$. At the two rational endpoints
$y_0\pm10^{-240}$ the truncated expression plus this
infinite tail is respectively strictly positive and strictly
negative. The code records all twenty-two strict signs.
Continuity produces a zero within each bracket. Every zero
of $\wp$ is one of the pair $\pm P$, and $N$ agrees on
that pair, so no unproved choice of root is needed.

\item The product formula on this bracket has positive real factors:
\begin{align*}
 \log N={}&\log(2\cosh\pi y)-\frac{\pi T}{6}
                      -\frac{\pi y^2}{T}\\
 &+\sum_{n\ge1}\{\log(1+q^ne^{2\pi y})
                         +\log(1+q^ne^{-2\pi y})\}.
\end{align*}
Since $|\log(1+x)|\le |x|/(1-|x|)$ for real $|x|<1$,
the omitted tail is bounded by
\[
 \frac{2e^{2\pi Y}Q^{M+1}}
 {(1-Q)(1-e^{2\pi Y}Q^{M+1})}<10^{-1260}.
\]
All calculations use Arb enclosures at $1152$ bits; the
radius comparisons and endpoint signs are exact interval
comparisons. The resulting radius is less than $10^{-236}$.
Adding the exact finite local logarithms of \ref{ari:C.6.3} and
subtracting this enclosure gives the claimed height bound.

\item For completeness, the singular-modulus certificate uses
$E_4=1+240\sum n^3q^n/(1-q^n)$ and
$\Delta_0=q\prod(1-q^n)^{24}$. Their respective error
bounds, after truncation at $M$, are
\[
 \frac{240(M+1)^3Q^{M+1}(1+4Q+Q^2)}{(1-Q)^5},
 \qquad
 \frac{24Q^{M+1}}{(1-Q)(1-Q^{M+1})}
\]
for $E_4$ and for the logarithm of the product. The first
uses $(M+1+k)^3\le(M+1)^3(k+1)^3$ and the differentiated
geometric series; the second uses the same logarithm bound.
Thus $E_4^3/\Delta_0$ is rigorously enclosed. The complete
certificate is \path{code/C_cm_height_arb.py}, with its
output and rational bracket centres retained.
\end{enumerate}
\end{proof}


\subsection{Supplement: ranks, saturation, and central derivatives}

We retain the order of the eleven rational CM values in the stated conventions, and
write $E_i$ for the global minimal models recorded in the preceding
tables. The real period in this subsection is
\[
 \Omega_i=\int_{E_i(\R)}|\omega_i|,
\]
including both connected components when $E_i(\R)$ is disconnected.
The canonical height $\widehat h$ is relative to $(O)$; PARI's
\texttt{ellheight} and the regulator convention used in the BSD formula
are relative to $2(O)$, and hence equal $2\widehat h$ in rank one.
This factor of two is essential in every comparison below.

\begin{proposition}[ARI \ref{ari:C.6.R1}; Proved by descent and full saturation]\label{ari:C.6.R1}
\StatementRecord{ari:C.6.R1}{proved}{Historical events: GPT-6 Astra (OpenAI), ChatGPT Pro mode, effort Pro, S6\_09\_PAPERS, 17 November 2026: targeted proof check. Also checked in this writing. Pre-continuation writing events: GPT-6 Astra (OpenAI), ChatGPT Pro mode, effort Pro, S6\_10\_PAPERS, 18 November 2026: computation or formal certificate; targeted proof check. Continuation metadata events: GPT-6 Astra (OpenAI), ChatGPT Pro mode, effort Pro, S6\_10\_PAPERS continuation, 19 November 2026: used and rechecked in this run; status-vocabulary reconciliation; metadata reconciliation only; no new mathematical check. \emph{Pre-continuation recorded proof scope.} Read complete printed arithmetic proofs and re-derived the exact scaling/local-height equations and the genus-trace square denominator argument. Rank, saturation, directed L-derivative replays and precise external theorem matches remain explicit gates; no proof verdict is awarded to those computational claims by this read. Reran the exact model/local-height programs in PARI2.17.2 and the unchanged eclib20250122 saturation code. All eleven displayed models and local coefficients match; ranks are exactly2 for rows1,8 and exactly1 for the remaining nine. All nine saturation indices are1, with no remaining primes. The full actual outputs and source hashes are retained. Opened PARI ellrank/ellheight normalization documentation and Cremona--Siksek Theorem1, author-PDFp8, for the algorithmic scope.}{the hypotheses, notation and definitions stated here; \citet[Theorem 1]{CrCremonaSiksek2006}}

The algebraic and analytic ranks agree for the eleven curves. They are
two for $j=8000$ and $j=-12288000$, and one for all nine other values.
For every rank-one curve the section $P$ is a primitive generator of
$E(\Q)/E(\Q)_{\rm tors}$. Its index is therefore one. For the two
rank-two curves the index of $\Z P$ in this free quotient is infinite.
\[
\begin{array}{r|r|r|r}
i&j&\operatorname{rank}E(\Q)&[E(\Q)/E(\Q)_{\rm tors}:\Z P]\\\hline
1&8000&2&\infty\\
2&54000&1&1\\
3&287496&1&1\\
4&16581375&1&1\\
5&-3375&1&1\\
6&-32768&1&1\\
7&-884736&1&1\\
8&-12288000&2&\infty\\
9&-884736000&1&1\\
10&-147197952000&1&1\\
11&-262537412640768000&1&1
\end{array}
\]
\end{proposition}


\begin{proof}[Arithmetic verification]
\begin{enumerate}[label=\arabic*.,leftmargin=*]
\item Exact $2$-descent in PARI/GP returns equal lower and upper rank bounds
for all eleven curves; its entire output, including independent points,
is in \path{code/C_cm_rank_lvalues.out}. Independent eclib descent
and saturation were completed for the first nine curves. The separate
full-saturation calculation \path{code/Crank_saturation.cpp} treats
all nine rank-one points, using eclib's automatic height-based index
bound and checking every prime up to that bound. All nine returned
index one and an empty list of unsaturated primes.

\item For clarity, the rounded-down bounds for the indices in the subgroup
of points having good reduction everywhere are respectively
\[
 1,1,1,2,2,2,3,5,5
\]
in the rank-one order $i=2,3,4,5,6,7,9,10,11$; the algorithm also
checks the primes dividing component-group orders. The underlying
height lower bound is the one in
\citet[Theorem 1]{CrCremonaSiksek2006}, applied to a minimal
Weierstrass equation and the subgroup
$E(\Q)\cap E^0(\R)\cap\bigcap_pE^0(\Q_p)$.
For the largest curve the lower bound on the nonzero height in this
subgroup is $0.651940632595967949\ldots$, its regulator is
$17.323913974368488616\ldots$, and the resulting index bound is
$5.154885040175890631\ldots$. The finite-field saturation maps for
$2,3,5$ already have full rank at residue primes $17,17,59$.
These data certify primitivity; merely finding $P$ in the output of
\texttt{ellrank} would not do so. The complete height-bound and
saturation transcripts are retained, including the independent
descent run that was stopped after its ninth completed curve.

\item The analytic ranks follow from the explicit derivative sums and bounds
proved below. The two rank-two bases furnished by descent are
$(-7,441),(105,833)$ for $i=1$, and
$(0,32004),(22218,3305456)$ for $i=8$.
\end{enumerate}
\end{proof}

\begin{proposition}[ARI \ref{ari:C.6.R2}; Proved: central derivatives with explicit tails]\label{ari:C.6.R2}
\StatementRecord{ari:C.6.R2}{proved}{Historical events: GPT-6 Astra (OpenAI), ChatGPT Pro mode, effort Pro, S6\_09\_PAPERS, 17 November 2026: targeted proof check. Also checked in this writing. Pre-continuation writing events: GPT-6 Astra (OpenAI), ChatGPT Pro mode, effort Pro, S6\_10\_PAPERS, 18 November 2026: computation or formal certificate; targeted proof check. Continuation metadata events: GPT-6 Astra (OpenAI), ChatGPT Pro mode, effort Pro, S6\_10\_PAPERS continuation, 19 November 2026: used and rechecked in this run; status-vocabulary reconciliation; metadata reconciliation only; no new mathematical check. \emph{Pre-continuation recorded proof scope.} Read complete printed arithmetic proofs and re-derived the exact scaling/local-height equations and the genus-trace square denominator argument. Rank, saturation, directed L-derivative replays and precise external theorem matches remain explicit gates; no proof verdict is awarded to those computational claims by this read. Read the complete retained C++ summation program and replayed both segmented runs unchanged:384bits/radius128/cutoff170 and512bits/radius96/cutoff190. The latter reproduces915818441 coefficients,69802132 nonzero values,sumabs1904681267996,4745134 blocks and the saved derivative exactly. Explicit tail is3.09321800754876143967e-78. Independently enclosed all eleven real periods and fresh derivative ratios with Arb; all displayed fifty-digit values pass. Opened Rubin Proposition5.3 p15 and Theorem10.1/full Section10 pp47–50; the exact CM isogeny and central-vanishing application are now written out.}{the hypotheses, notation and definitions stated here; \cref{ari:eq:Cr-tails}; \citet[Section 3]{CrDarmon2009}; \citet[Theorem 10.1]{CrRubin1999}}

The following are the values of $L'(E_i,1)/\Omega_i$; zeros in rows
one and eight are exact. More digits, the periods, derivatives and
tail bounds are in \path{results_Crank.csv}.
\begin{center}\small
\begin{tabular}{r|l}
$i$&$L'(E_i,1)/\Omega_i$\\\hline
1&$0$\\
2&$5.134398425737889664345619204085016979705061236171249722654226\ldots$\\
3&$2.134120849044754020051650801945664635106456422166881627967598\ldots$\\
4&$31.81013979419927776012209427941545866492237285892013668350005\ldots$\\
5&$1.869462708725885120030080805051387134542481374813232701178304\ldots$\\
6&$3.428873379050567976208155749688133688667614469824389504923957\ldots$\\
7&$0.362583348071598548315763082767526599722077924892872839203426\ldots$\\
8&$0$\\
9&$9.629273918234131678053742152263035417231796757756973886621556\ldots$\\
10&$16.04238857830583792513538634637663343755958963581574670348088\ldots$\\
11&$1247.321806154531180413038860827393095408344267114181914400926\ldots$
\end{tabular}
\end{center}
The two second derivatives, divided by the same real periods, are
\begin{align*}
 L''(E_1,1)/\Omega_1
 &=28.71408722758099118128474692322079946598060512912402172882525\ldots,\\
 L''(E_8,1)/\Omega_8
 &=78.03222041637032223732443782051140882655135939512581188900823\ldots.
\end{align*}
In particular the leading Taylor coefficients for these two curves
are half the displayed second derivatives.
\end{proposition}


\begin{proof}[Derivative formulas and error control]

\begin{enumerate}[label=\arabic*.,leftmargin=*]
\item Let $N$ be the conductor, $c=2\pi/\sqrt N$, and
$L(E,s)=\sum_{n\geq1}a_nn^{-s}$. Put
\[
 E_1(x)=\int_x^\infty\frac{e^{-u}}u\,du,
 \qquad G(x)=\int_x^\infty e^{-u}\frac{\log(u/x)}u\,du.
\]
Splitting the Mellin integral at $1/\sqrt N$ and using the functional
equation gives, for root number $-1$,
\begin{equation}\label{ari:eq:Cr-first}
 L'(E,1)=2\sum_{n\geq1}\frac{a_n}{n}E_1(cn).
\end{equation}
For root number $+1$ and $L(E,1)=0$, it gives
\begin{equation}\label{ari:eq:Cr-second}
 L''(E,1)=4\sum_{n\geq1}\frac{a_n}{n}G(cn).
\end{equation}
One can verify the factors directly by writing the completed Mellin
integral as an integral over $[1,\infty)$ of
$(u^{s-1}+wu^{1-s})f(iu/\sqrt N)$: its first derivative has
the factor $2\log u$ when $w=-1$, and its second derivative
the factor $2\log^2u$ when $w=1$. Integration by parts gives
the two formulas above. The vanishing of the lower derivatives
removes all derivatives of the gamma and conductor factors.

\item Hasse's bound and multiplicativity imply
$|a_n|\leq d(n)\sqrt n\leq2n$. Also
$E_1(x)\leq e^{-x}/x$ and $G(x)\leq e^{-x}/x^2$.
The latter follows on writing $u=x+v$ and using
$\log(1+v/x)\leq v/x$. Consequently the absolute tails after
$n=M$ are at most
\begin{equation}\label{ari:eq:Cr-tails}
 \frac{4e^{-c(M+1)}}{c(M+1)(1-e^{-c})},\qquad
 \frac{8e^{-c(M+1)}}{c^2(M+1)^2(1-e^{-c})},
\end{equation}
respectively. For the first ten curves we used
$M=\lceil180/c\rceil$, exact integer coefficients, and 1024-bit
MPFR arithmetic. For the rank-two kernel we used the convergent series
\[
 G(x)=\frac12(\log x+\gamma)^2+\frac{\pi^2}{12}
       +\sum_{k\geq1}\frac{(-x)^k}{k^2k!}.
\]
After $k>x$ its terms decrease in absolute value. The alternating
remainder was cut below $10^{-140}$. At most 740 terms were needed;
the largest intermediate term was less than $1.4\cdot10^{72}$.
The 1024-bit arithmetic therefore leaves far more than 100 digits
after this cancellation. The largest bound in
\eqref{ari:eq:Cr-tails} was $3.46\cdot10^{-77}$. The source and complete
outputs are \path{code/Crank_coefficients.gp},
\path{code/Crank_derivatives.cpp}, and their corresponding outputs.

\item For $i=11$, a full vector of $819$ million coefficients would be
unnecessarily large. We instead evaluated the CM theta series in
bounded-memory segments. The curve is the quadratic twist, by
$d=185801=7\cdot11\cdot19\cdot127$, of the conductor-$163^2$
curve
\[
 y^2+y=x^3-2174420x+1234136692.
\]
Write $D=163$, $\chi_D(a)=(a/D)$, and $\chi_d(n)=(d/n)$.
The exact coefficient $a_n$ is the sum of
\[
 \chi_d(n)\chi_D(a)a \quad(n=a^2,\ a>0)
\]
and
\[
 \chi_d(n)\chi_D((D+1)u/2)u
 \quad\left(n=\frac{u^2+Dv^2}{4},\ u,v>0,
                           \ u\equiv v\pmod2\right).
\]
This is the Hecke-character theta series for
$\mathcal O_{\Q(\sqrt{-D})}$, grouping conjugate algebraic integers
and then applying the quadratic twist; see
\citet[Section 3]{CrDarmon2009}. The first 10000 coefficients agree
exactly with PARI for both the untwisted and twisted curve. The same
code at $D=67,d=217$ reproduces the independent PARI derivative.

\item To accelerate the sum, in a block $n=n_0+r$, $|r|\leq R$, form the
exact integer moments $M_k=\sum a_nr^k$. For $x=cn_0$ write
\[
 e_k=e^{-x}(-x)^k/k!,\quad g_0=e^{-x},\quad
 g_k=e_k-g_{k-1},\quad
 \phi_0=E_1(x),\quad\phi_k=-\phi_{k-1}-g_{k-1}/k.
\]
Then the block contribution to $L'/2$ is
$\sum_k\phi_k M_k/n_0^{k+1}$. This follows by Taylor expansion
of $E_1(x(1+s))/(1+s)$. On $|s|=1/2$ its modulus is at most
$2E_1(x/2)$, which is less than $14$ throughout the block range.
We take enough terms that $(2R/n_0)^{k+1}<10^{-85}$.
The Cauchy remainder, summed with the actual
$\sum|a_n|<2\cdot10^{12}$ and $n_0>7700$, is less than
$10^{-75}$. The values of $E_1$ at consecutive block centres
are continued by the analogous Taylor expansion and refreshed
from a correctly rounded MPFR evaluation every 1024 blocks.
The continuation bound is $10^{-110}$ per normalized expansion.
All moments are integers, not floating-point sums.

\item The 384-bit run used $R=128$ and cutoff $170$; a separate 512-bit
run used $R=96$ and cutoff $190$. Their values agree beyond 75
decimal places. The latter used $915818441$ coefficients in
segments, with an omitted tail less than $3.10\cdot10^{-78}$,
and gave
\[
 L'(E_{11},1)=
 2.296582198163380608965915774436446055185027904705664673866754\ldots.
\]
The elementary operation-rounding bounds at these precisions are
smaller than $10^{-90}$: in the recurrences above
$|e_k|\leq1$, $|g_k|\leq1$, and $|\phi_k|<20$;
there are fewer than $5\cdot10^6$ blocks and the recurrences are
refreshed before 1024 consecutive continuations. Bounding each
normalized moment by $\sum|a_n|(R/n_0)^k$ gives this estimate
without assuming cancellation in the total sum. The arithmetic
and truncation margins are amply sufficient for every displayed
50-digit value. Both executable builds and both transcripts are
retained as \path{code/Crank_cm_stream*}.

\item The real periods are enclosed from the exact minimal invariants,
independently of the saved decimal period values.  For rows $1,2,3,4$
put respectively $T=\sqrt2,\sqrt3,2,\sqrt7$ and $\tau=iT$;
for rows $5,\ldots,11$ put
$T=\sqrt{7}/2,\sqrt{11}/2,\sqrt{19}/2,\sqrt{27}/2,
\sqrt{43}/2,\sqrt{67}/2,\sqrt{163}/2$ and $\tau=1/2+iT$.
Let $r=e^{-2\pi T}$ and $q=e^{2\pi i\tau}$.  The identity
$E_4=1+240\sum_{n\geq1}\sigma_3(n)q^n$ and
$\sigma_3(n)<2n^3$ give, after $M=100$ terms, the tail bound
\[
480r^{A}\left\{
\frac{A^3}{1-r}+\frac{3A^2r}{(1-r)^2}
+\frac{3Ar(1+r)}{(1-r)^3}
+\frac{r(1+4r+r^2)}{(1-r)^4}\right\},\qquad A=101.
\]
This follows by expanding $(A+j)^3$ and summing each geometric
moment.  All eleven parameters satisfy
$1-1008r/(1-32r)>0$; the bound
$\sigma_5(n)<2n^5$ and $n^5\leq32^{n-1}$ therefore proves
$E_6(\tau)>0$.  The exact Weierstrass invariant relations yield
$\kappa=2\pi(E_4/c_4)^{1/4}$.  For $c_6>0$ the real lattice
generator has length $\kappa$.  For $c_6<0$ the lattice is rotated
by $i$: its real generator has length $T\kappa$ in the rectangular
case and $2T\kappa$ in the other case.  Indeed the real vectors in
$i(\Z+iT\Z)$ are $T\Z$, while those in
$i(\Z+(1/2+iT)\Z)$ are $2T\Z$.  Multiplication by the number
$n_{\R}=2$ for $\Delta>0$ and $n_{\R}=1$ for $\Delta<0$
gives the full real period used in the statement.

The current program \path{code/I/certify_all_period_ratios.py}
performs these operations with 768-bit Arb intervals and verifies
every displayed period and derivative ratio.  For the first ten
derivatives it enlarges the explicit tail by one percent and adds
$10^{-100}$ for arithmetic and kernel-series errors.  This last
bound is conservative: fewer than $5\cdot10^5$ coefficients and
740 series terms occur, the intermediate terms are below
$1.4\cdot10^{72}$, and the 1024-bit accumulated rounding and
$10^{-140}$ kernel remainders are below this allowance.
For the last derivative it adds $1.1\cdot10^{-75}$, covering the
proved block Taylor and rounding bounds above.  Division by the
strictly positive period intervals preserves all fifty displayed
digits.  The resulting intervals and exact invariants are in
\path{code/I/records/all_period_ratios_arb.json}.

\item The exact root numbers are $+1$ for $i=1,8$ and $-1$ otherwise.
All points are nontorsion by Statement C.5. For a CM curve,
the Coates--Wiles theorem, in the form
\citet[Theorem 10.1]{CrRubin1999}, says that a nonzero central
Hecke $L$-value forces the Mordell--Weil group over the CM field
to be finite. Proposition~5.3 of the same source supplies, for a nonmaximal CM order, an isogenous
curve with maximal order: the quotient is by the Galois-stable subgroup attached to the conductor ideal. The isogeny induces an isomorphism on Mordell--Weil groups tensored with $\Q$ and on their rational Tate modules, preserving ranks and local $L$-factors.
Since $L(E/K,s)=L(E/\Q,s)^2$ for the CM field $K$, nontorsion
therefore forces $L(E,1)=0$. This proves the needed central
vanishing also in the two sign-$+1$ cases. The functional equation
then forces $L'(E,1)=0$ for those two curves. The positive first
or second derivatives just evaluated prove exactly the asserted
analytic ranks.
\end{enumerate}

\end{proof}

\begin{proposition}[ARI \ref{ari:C.6.R3}; Proved: the consequences of Gross--Zagier and Kolyvagin]\label{ari:C.6.R3}
\StatementRecord{ari:C.6.R3}{proved}{Historical events: GPT-6 Astra (OpenAI), ChatGPT Pro mode, effort Pro, S6\_09\_PAPERS, 17 November 2026: targeted proof check. Also checked in this writing. Pre-continuation writing events: GPT-6 Astra (OpenAI), ChatGPT Pro mode, effort Pro, S6\_10\_PAPERS, 18 November 2026: targeted proof check. Continuation metadata events: GPT-6 Astra (OpenAI), ChatGPT Pro mode, effort Pro, S6\_10\_PAPERS continuation, 19 November 2026: used and rechecked in this run; status-vocabulary reconciliation; metadata reconciliation only; no new mathematical check. \emph{Pre-continuation recorded proof scope.} Read complete printed arithmetic proofs and re-derived the exact scaling/local-height equations and the genus-trace square denominator argument. Rank, saturation, directed L-derivative replays and precise external theorem matches remain explicit gates; no proof verdict is awarded to those computational claims by this read. Matched Gross–Zagier I.7.3 pp11–12 and final proof p129 visually, Gross Theorem1.3 pp235–236 visually, and Jetchev–Lauter–Stein Theorem1.1 p2. Re-derived the rational differential/real-component period factor and added the explicit restriction–inflation argument carrying finite Sha from the Heegner field to Q. Exact derivative and algebraic-rank inputs were freshly replayed. The claimed nine positive rational factors and finite Sha conclusions are supported at this complete receiving-proof scope.}{the hypotheses, notation and definitions stated here; \citep{CrGrossZagier1986}; \citet[Theorem 1.3]{CrGross1991}; \citet[Theorem 1.1]{CrJetchev2008}}

For each of the nine rank-one curves there is a positive rational
number $q_i$ such that
\begin{equation}\label{ari:eq:Cr-rational-factor}
 \frac{L'(E_i,1)}{\Omega_i}=q_i\widehat h(P_i).
\end{equation}
Their Tate--Shafarevich groups over $\Q$ are finite. Consequently
the corresponding $D_{\rm part}$ is an explicit logarithm minus
$(6/(\pi q_i))L'(E_i,1)/\Omega_i$. For the two rank-two curves, $L'(E_i,1)=0$ while
$\widehat h(P_i)>0$.
\end{proposition}


\begin{proof}

\begin{enumerate}[label=\arabic*.,leftmargin=*]
\item Gross--Zagier's theorem applies because each curve is modular and
$L(E_i,1)=0$. CM modularity follows from its Hecke-character
theta series. Their Theorem I.7.3, in the opened original
preprint \citep{CrGrossZagier1986}, gives a rational point $R$ and
a nonzero rational factor relating $L'(E,1)$ to its canonical
height times the real period. The differentials in the two conventions are rational multiples:
write $\omega'=a\omega_E$ with $a\in\Q^*$.  A primitive real
period for $\omega'$ is $|a|\Omega_i/n_{\R}$, where
$n_{\R}\in\{1,2\}$ is the number of real components.
Thus the change to the full real period multiplies the rational
factor by $|a|/n_{\R}$.
As $L'(E_i,1)>0$, $R$ has infinite order. Primitivity gives
$R=mP_i+T$, $m\ne0$, $T$ torsion; its height is
$m^2\widehat h(P_i)$. This proves
\eqref{ari:eq:Cr-rational-factor}, including positivity.

\item The rank-one Gross--Zagier--Kolyvagin implication asserts rank one
and finite Tate--Shafarevich group when the central zero is simple.
One can use Kolyvagin's Heegner-point theorem as stated in
\citet[Theorem 1.3]{CrGross1991}, together with the nonzero
Gross--Zagier Heegner height; a modern primary research statement
of the rank conclusion is \citet[Theorem 1.1]{CrJetchev2008}.
The hypotheses are satisfied separately by the positive derivative
and exact sign in each of the nine rows.  Explicitly, the auxiliary
imaginary quadratic field $K$ in the proof of Gross--Zagier I.7.3
has split Heegner primes and $L(E_i^D,1)\ne0$; hence
$L'(E_i/K,1)=L'(E_i,1)L(E_i^D,1)\ne0$.  Its Heegner point has
positive height, so Gross's Theorem~1.3 gives finite
$\operatorname{Sha}(E_i/K)$.  The restriction map
\[
\operatorname{Sha}(E_i/\Q)\longrightarrow
\operatorname{Sha}(E_i/K)
\]
has kernel contained in
$H^1(\operatorname{Gal}(K/\Q),E_i(K))$ by restriction--inflation.
This last group is finite: the Mordell--Weil group is finitely
generated, and its degree-one cohomology for this order-two group
is finitely generated and killed by two.  The finite kernel and
finite image prove the claimed finiteness over $\Q$.
For the two other rows the established equations are
$L'(E_i,1)=0$ and $\widehat h_O(P_i)>0$; therefore an equation
$L'(E_i,1)/\Omega_i=q\widehat h_O(P_i)$ with $q>0$ would contradict
these exact values.  Their rank-two regulator, for any free basis
$P,Q$, is the determinant
$\widehat h_O(P)\widehat h_O(Q)-\langle P,Q\rangle^2$.
\end{enumerate}

\end{proof}

\paragraph{What the numerical BSD comparison does and does not determine.}
For a primitive rank-one point, the BSD prediction in these conventions is
\[
 \frac{L'(E,1)}{\Omega_E}
 =\frac{2\#\operatorname{Sha}(E/\Q)\prod_pc_p}
        {\#E(\Q)_{\rm tors}^{\,2}}\,\widehat h(P).
\]
The independently evaluated quotient after dividing by
$2\widehat h(P)\prod_pc_p/\#E(\Q)_{\rm tors}^2$ is one, to
more than 69 digits, for $i=2,3,4,5,6,7,9,10$, and nine for
$i=11$. Thus the numerically predicted rational factors $q_i$ are
\[
 8,4,32,4,8,4,16,16,576,
\]
in that order. This initial numerical comparison alone proves neither
these exact factors nor the orders $1$ or $9$ of $\operatorname{Sha}$.
The exact factors are subsequently proved in \ref{ari:C.6.R4} and \ref{ari:C.6.R5}:
the latter supplies an independent Gross--Zagier denominator bound
and rigorous intervals for all remaining ratios. This determines
the analytic BSD quotients exactly, while the orders of
$\operatorname{Sha}$ are proved here only in the three cases of
\ref{ari:C.6.R4}. In particular primitivity alone does not determine
$\#\operatorname{Sha}$ for the last curve.

\begin{cited}[ARI \ref{ari:C.6.R4}; Known input: three exact leading factors]\label{ari:C.6.R4}
\StatementRecord{ari:C.6.R4}{proved}{Historical events: GPT-6 Astra (OpenAI), ChatGPT Pro mode, effort Pro, S6\_09\_PAPERS, 17 November 2026: targeted proof check. Also checked in this writing. Pre-continuation writing events: GPT-6 Astra (OpenAI), ChatGPT Pro mode, effort Pro, S6\_10\_PAPERS, 18 November 2026: targeted proof check; used and rechecked in this run. Continuation metadata events: GPT-6 Astra (OpenAI), ChatGPT Pro mode, effort Pro, S6\_10\_PAPERS continuation, 19 November 2026: used and rechecked in this run; status-vocabulary reconciliation; metadata reconciliation only; no new mathematical check. \emph{Pre-continuation recorded proof scope.} Read complete printed arithmetic proofs and re-derived the exact scaling/local-height equations and the genus-trace square denominator argument. Rank, saturation, directed L-derivative replays and precise external theorem matches remain explicit gates; no proof verdict is awarded to those computational claims by this read. Read Miller Definition1.1/Theorem1.2 p2, Section7 opening p16 and complete exception Table9 p23. Conductors3136,3136,361 lie outside the exceptions; all have proved analytic rank1. Exact local/torsion data give coefficient4. Current Arb height and derivative/period intervals isolate the resulting positive integer Sha order as1 in all three rows. Expanded known application remains in the unchanged proof compendium.}{the exact models and primitivity in \cref{ari:C.6.R1};
the certified derivatives in \cref{ari:C.6.R2};
\citet[Definition~1.1, Theorem~1.2, and Table~9]{miller2011};
the full application retained in the proof compendium, Paper~I,
label \nolinkurl{app:rederive:ari-C.6.R4}}

For $i=3,5,7$, namely $j=287496,-3375,-884736$, one has
\[
 \operatorname{Sha}(E_i/\Q)=0,\qquad
 \frac{L'(E_i,1)}{\Omega_i}=4\widehat h_O(P_i).
\]
Thus the factors $q_3=q_5=q_7=4$ are determined exactly by the integer bounds in the cited application. For the smallest of these examples,
$E:y^2+y=x^3-38x+90$ and $P=(0,9)$, this gives the exact
reduction
\[
 N(P)=\exp\left(-\frac{L'(E,1)}{4\Omega_E}\right).
\]
\end{cited}

The expanded source derivation is retained in the proof compendium, Paper~I, label \nolinkurl{ari:C.6.R4}.




The three exact cases specialize \citet[Definition~1.1, Theorem~1.2 and Table~9]{miller2011}. The complete application is retained in the proof compendium, Paper~I, label \nolinkurl{app:rederive:ari-C.6.R4}.


\paragraph{The scope of Rubin's theorem.}
The applicable ordinary CM results concern a prime $p$ at which
$E$ has good ordinary reduction, equivalently $p$ splits in the CM
field. One precise opened formulation of Rubin's rank-one formula
is \citet[Theorem 1]{CrBDP2012}: for a curve over $\Q$ with CM
by the maximal order, an ordinary prime prime to the conductor,
and a simple central zero, an appropriate Katz $p$-adic $L$-value
is, up to a nonzero algebraic factor, the square of a formal
logarithm of a nontorsion rational point divided by a $p$-adic
period. Its input is analytic rank one, already verified here.
It is a $p$-adic identity; it does not identify the real exponential
$\exp\widehat h(P)$ as algebraic.

For the $p$-part of the real leading-term formula, a modern precise
version is \citet[Theorem A]{CrCastella2025}. In addition to analytic
rank one it requires full CM, abelian torsion field over the CM
field, $p\nmid6h_K$, splitting of $p$, conductor prime to $p$,
and that the different of $K$ divide the Hecke conductor exactly.
For the maximal odd-discriminant cases here, the latter follows
from the exponent two of the ramified prime in the rational
conductor. The nonmaximal cases of discriminant $-12$ and $-28$
can be passed through a $2$-isogeny at $p>3$. The discriminant
$-16$ case requires care: the direct application of this version
has a different-conductor hypothesis at $2$ that fails, and we
do not claim that this particular formulation proves its full
leading-term formula. Rubin's original ordinary main-conjecture
paper is \citep{CrRubin1991}; its bibliographic record was checked,
but its complete theorem text was not available in the opened
publisher copy, so no additional hypothesis-free claim is based
on it here.

In particular these ordinary-prime statements do not prove that
the $3$-primary part of $\operatorname{Sha}$ has order $9$ for the
last curve: $3$ is inert in $\Q(\sqrt{-163})$ and also excluded
from the stated $p>3$ hypotheses. \ref{ari:C.6.R5} proves the corresponding
analytic quotient to be exactly $9$ by a separate denominator
argument. The identification with the actual Tate--Shafarevich
group remains a separate arithmetic question in that row.


\subsection{Supplement: exact leading factors from a genus-character denominator bound}

We retain the full real period $\Omega_i$ and the absolute canonical
height $\widehat h_O$ of the preceding subsection. The argument below
uses the Gross--Zagier formula to bound a rational denominator; it does
not assume the leading-term Birch--Swinnerton-Dyer formula.

\begin{proposition}[ARI \ref{ari:C.6.R5}; Proved: all remaining rational leading factors]\label{ari:C.6.R5}
\StatementRecord{ari:C.6.R5}{proved}{Historical events: GPT-6 Astra (OpenAI), ChatGPT Pro mode, effort Pro, S6\_09\_PAPERS, 17 November 2026: targeted proof check. Also checked in this writing. Pre-continuation writing events: GPT-6 Astra (OpenAI), ChatGPT Pro mode, effort Pro, S6\_10\_PAPERS, 18 November 2026: computation or formal certificate; targeted proof check; used and rechecked in this run. Continuation metadata events: GPT-6 Astra (OpenAI), ChatGPT Pro mode, effort Pro, S6\_10\_PAPERS continuation, 19 November 2026: used and rechecked in this run; status-vocabulary reconciliation; metadata reconciliation only; no new mathematical check. \emph{Pre-continuation recorded proof scope.} Read complete printed arithmetic proofs and re-derived the exact scaling/local-height equations and the genus-trace square denominator argument. Rank, saturation, directed L-derivative replays and precise external theorem matches remain explicit gates; no proof verdict is awarded to those computational claims by this read. Reran B\_C6\_gz\_exact\_metadata.gp with only the unavailable database-label display omitted. All six exact twist/differential/conductor/coprimality/Heegner-splitting identities pass. All four CM theta coefficient vectors and auxiliary quadratic-twist vectors agree with independently computed elliptic Fourier coefficients through n=10000, beyond each printed Sturm bound. Directed numerical enclosures and the remaining primary theorem normalizations are still being checked. All six current directed MPFR derivative/auxiliary central-value enclosures and Arb exact-factor recognitions passed. Ratios r=64,64,8,16,400,64 and q=8,32,8,16,16,576; genus-trace magnitudes32,64,8,16,80,192. Fresh GP Fourier coefficients and exact eleven-model comparison passed. An initial auxiliary script used PARI reserved variable I and failed; changed to idx and reran all five affected file certificates successfully. No failed output was accepted. Opened original Cai-Shu-Tian2014 Theorem1.1 pp2524-2525 and Belabas-Bernardi-Perrin-Riou2019 AppendixB p102 unnumbered26569a1/a2 block. Verified its Manin constant1 and modular degree59368 visually. Height-normalization and remaining modular-input checks remain separately tracked. Recomputed all first-five small-base Manin constants by installed PARI2.17.2 ellmaninconstant(E); all equal1. Local runtime help explicitly states unconditional algorithm. Initial two-argument call from newer web docs was rejected; preserved failed logs, removed unsupported flag and reran five computations successfully. Opened author-uploaded YZZ2011 Section7.1.1 printedpp218-219 including Theorems7.1.1/7.1.2 and Proposition7.1.3; confirmed base-field-relative height and pairing without half, yielding exact factor4 relative to absolute O-height. Opened Darmon2009 Section3pp4-5 defining canonical CM Hecke theta series. Opened Stein Theorem9.18 and its entire proof pp171–173, Corollaries9.19–9.20 pp173–174, matching the coefficient bounds and coset-product argument. Together with the separately recorded YZZ, Cai–Shu–Tian, Darmon, Manin-constant and six directed numerical checks, this closes the named primary-source gates for the six denominator recognitions.}{the hypotheses, notation and definitions stated here; \cref{ari:C.6.R4}; \citet[Appendix B, p.~102, unnumbered $26569\mathrm a1/26569\mathrm a2$
block]{CgzBelabas2019}; \citet[Theorem 1.1, pp.~2524--2525]{CgzCai2014}; \cref{ari:C.6.R2,ari:C.6.R3}; \cref{ari:C.6.R1}; \citet[\S7.1.1, pp.~218--219, Theorem 7.1.2]{CgzYZZ2011}; \cref{ari:eq:Cgz-product}; \citet[\S3, pp.~4--5]{CrDarmon2009}; \citet[Theorem 9.18]{CgzStein2007}; \cref{ari:eq:Cr-first}; \cref{ari:C.6.5}; \cref{ari:eq:Cgz-central-integrality}; \cref{ari:eq:Cgz-denominator}}

For the six rank-one rows $i=2,4,6,9,10,11$, one has
\[
 \frac{L'(E_i,1)}{\Omega_i}
   =q_i\widehat h_O(P_i),\qquad
 (q_2,q_4,q_6,q_9,q_{10},q_{11})=(8,32,8,16,16,576).
\]
Together with \ref{ari:C.6.R4}, this determines every rank-one factor:
\[
\begin{array}{c|rrrrrrrrr}
 i&2&3&4&5&6&7&9&10&11\\\hline
 q_i&8&4&32&4&8&4&16&16&576.
\end{array}
\]
The certified auxiliary ratios and denominator bounds used to identify
the six additional factors are as follows. Here $D$ is the auxiliary
negative fundamental discriminant,
$r_i=L(E_i^D,1)/\Omega_{i,D}$, and the denominator of $q_i$ is at most
$B_i$.
\[
\begin{array}{c|r|r|r|r|r}
 i&D&r_i&B_i&q_i&|a_i|\\\hline
 2&-23&64&256&8&32\\
 4&-31&64&256&32&64\\
 6&-19&8&32&8&8\\
 9&-19&16&64&16&16\\
 10&-3&400&1600&16&80\\
 11&-3&64&256&576&192
\end{array}
\]
The last column is the absolute genus-trace integer in the proof. These
identities determine the analytic BSD quotients exactly: they are one
for $i=2,4,6,9,10$ and nine for $i=11$. They do not by themselves
identify those numbers with the orders of the Tate--Shafarevich groups.
\end{proposition}



\begin{proof}
\begin{enumerate}[label=\arabic*.,leftmargin=*]
\item We give the arithmetic denominator bound first and then the interval
certificate that isolates the listed rational numbers.

\item \paragraph{1. A parametrization under quadratic twisting.}
Let $E_0/\Q$ have conductor $N_0$, and suppose a parametrization
$\phi:X_0(N_0)\to E_0$ sends $\infty$ to $O$ and satisfies
\[
 \phi^*\omega_0=c_0f_0(q)\,dq/q,
\]
where $\omega_0$ is a minimal differential and $f_0$ is the normalized
newform. Let $b$ be an odd fundamental discriminant coprime to $N_0$,
and suppose that the minimal twist has conductor $N_0b^2$ and exact
invariants
\[
 c_4(E_0^b)=b^2c_4(E_0),\qquad
 c_6(E_0^b)=b^3c_6(E_0).
\]
Then $E_0^b$ has a parametrization from $X_0(N_0b^2)$ with the same
Manin factor $c_0$.

\item Indeed, put $m=|b|$ and, using addition on $E_0$, define
\[
 F(z)=\sum_{a\bmod m}\chi_b(a)\phi(z+a/m).
\]
For $\gamma=\left(\begin{smallmatrix}A&B\\C&D_0\end{smallmatrix}\right)$
in $\Gamma_0(N_0m^2)$ choose $a'\equiv aD_0^2\pmod m$.
Direct multiplication gives
$T_{a/m}\gamma T_{a'/m}^{-1}\in\Gamma_0(N_0)$, where $T_t$
denotes translation by $t$. Because $\chi_b(D_0^2)=1$, the sum is
invariant. Each summand is an algebraic morphism on a common
congruence cover; invariance therefore descends the sum to
$X_0(N_0m^2)$ over $\Q(\zeta_m)$. Its expansion at $\infty$
gives $F^{\sigma_u}=\chi_b(u)F$ for
$\sigma_u(\zeta_m)=\zeta_m^u$. Quadratic descent gives the required
map to $E_0^b$ over $\Q$, with $\infty$ mapping to $O$.
On differentials, the Gauss-sum identity gives
\[
 F^*\omega_0=c_0\tau(\chi_b)f_b(q)\,dq/q,
 \qquad \tau(\chi_b)^2=b.
\]
For the twist isomorphism $\iota:E_0^b\to E_0$ the exact invariant
identities imply
$\iota^*\omega_0=\sqrt b\,\omega_b$ up to sign, with no additional
minimal-model scaling. Cancellation proves the assertion. The map
need not be optimal; optimality is not used below.

\item If $c_0=1$, let $t_b=\#E_0^b(\Q)_{\rm tors}$ and
$n_b=\#\pi_0(E_0^b(\R))$. Manin--Drinfeld says that the image of
the rational cuspidal divisor $[0]-[\infty]$ is torsion. Its elliptic
logarithm, along the imaginary axis, is $\pm L(E_0^b,1)$.
Consequently $t_bL(E_0^b,1)$ is an integral multiple of the positive
real lattice period $\Omega_b^+$, and
\begin{equation}\label{ari:eq:Cgz-central-integrality}
 t_bn_b\,\frac{L(E_0^b,1)}{\Omega_b}\in\Z,
 \qquad \Omega_b=n_b\Omega_b^+.
\end{equation}

\item \paragraph{2. The six small base curves and Heegner hypotheses.}
Use the following data; $E_i=E_0^d$ and $K=\Q(\sqrt D)$.
\[
\begin{array}{c|c|r|r|r}
 i&E_0&N_0&d&D\\\hline
 2&576\mathrm e2&576&-11&-23\\
 4&3136\mathrm r2&3136&57&-31\\
 6&121\mathrm b1&121&21&-19\\
 9&1849\mathrm a1&1849&21&-19\\
 10&4489\mathrm a1&4489&217&-3\\
 11&26569\mathrm a1&26569&185801&-3
\end{array}
\]
All six base Manin factors are one. For the first five bases this
is independently returned by PARI's unconditional modular-symbol
algorithm. For $26569\mathrm a1$, the value one and modular degree
$59368$ are explicitly tabulated by
\citet[Appendix B, p.~102, unnumbered $26569\mathrm a1/26569\mathrm a2$
block]{CgzBelabas2019}; their Algorithm 1 describes the computation.
The first five completed computations are in
\path{code/invariants_arithmetic/arithmetic/outputs/writing_B_C6_gz_exact_metadata.gp.txt}. The separate exact audit
\path{code/B_C6_gz_exact_metadata.gp} checks the minimal equations,
both $c_4,c_6$ scaling identities for $b=d,dD$, both conductor
identities, all torsion orders, and all splitting conditions. It
does not recognize any of these identities from floating-point
scaling factors.

\item In every row $d,D,dD$ are fundamental discriminants,
$\gcd(d,D)=\gcd(dD,N_0)=1$, and every prime dividing $N_0$ splits
in $K$. The genus character for $K(\sqrt d)/K$ factors through
$\operatorname{Pic}(\mathcal O_{|d|})$ and has conductor exactly
$|d|$. On ideals prime to $d$ it is
\[
 \chi(\mathfrak a)=\chi_d(N\mathfrak a).
\]
For a unit of the order locally at a prime $p\mid d$, its norm is
a square modulo $p$, so this character is trivial on the order's
units. It is nontrivial on the full local unit group: at a split
prime the norm is multiplication on two copies of $\Q_p^\times$,
and at an inert prime the residue-field norm onto $\mathbb F_p^\times$
is surjective. Since $p$ is odd and unramified in $K$, this also
proves that the conductor contains every $p\mid d$ exactly once.
Thus the primitive ring-class-character hypotheses of
\citet[Theorem 1.1, pp.~2524--2525]{CgzCai2014} hold. Their
$\mu(N_0,D)=0$ and
$u=[\mathcal O_{|d|}^{\times}:\Z^{\times}]=1$; the latter remains
true for $D=-3$ because $|d|>1$.

\item Let $P_{\mathfrak a}$ be the corresponding Heegner points and put
\[
 Y=\sum_{[\mathfrak a]}\chi(\mathfrak a)\phi(P_{\mathfrak a}).
\]
The coefficients are $\pm1$, so this is an actual algebraic point,
not merely an element of a tensor product with $\Q$. It is fixed
by $\ker\chi$ and anti-invariant under the other automorphism of
$K(\sqrt d)/K$. Twisting therefore gives $Y_d\in E_i(K)$.
Artin factorization gives
\[
 L(s,E_0,\chi)=L(E_i,s)L(E_i^D,s).
\]
The intervals below prove $L(E_i^D,1)>0$, hence the auxiliary
curve has rank zero. The rank-zero implication used here is the
same unconditional modular/CM rank implication used in \cref{ari:C.6.R2,ari:C.6.R3}.
Thus $E_i(K)$ has rank one. \ref{ari:C.6.R1} proves that $P_i$ is primitive
over $\Q$. Since $Y_d-\overline{Y_d}$ is torsion and
$Y_d+\overline{Y_d}$ is rational, there is an integer $a_i$ with
\[
 2Y_d=a_iP_i\pmod{E_i(K)_{\rm tors}},\qquad
 \widehat h_O(Y_d)=\frac{a_i^2}{4}\widehat h_O(P_i).
\]

\item \paragraph{3. The exact height and period normalization.}
Write $V_0=\operatorname{area}(\C/\Lambda_0)$ for the period area
of the minimal differential. Integrating $\phi^*\omega_0$ gives
\[
 4\pi^2(f_0,f_0)_{\Gamma_0(N_0)}
   =\frac{\deg\phi}{c_0^2}V_0.
\]
The degree in the denominator of CST's Theorem 1.1 cancels, giving
\begin{equation}\label{ari:eq:Cgz-product}
 L'(E_i,1)L(E_i^D,1)
  =\frac{2V_0}{c_0^2|d|\sqrt{|D|}}\,\widehat h_K(Y).
\end{equation}
Here $\widehat h_K$ is the diagonal N\'eron--Tate pairing relative
to $K$. This convention is explicit in
\citet[\S7.1.1, pp.~218--219, Theorem 7.1.2]{CgzYZZ2011}:
their underlying height is normalized by $[L:K]^{-1}$ over a
field of definition $L$, and their pairing is
$\widehat h_{\mathcal L}(x+y)-\widehat h_{\mathcal L}(x)
-\widehat h_{\mathcal L}(y)$.

\item One can also see the factor geometrically. For
$\mathcal L=\mathcal O(O)$, pull the normalized Poincar\'e bundle
back under $1\times\lambda_{\mathcal L}$. The result is
$m^*\mathcal L\otimes p_1^*\mathcal L^{-1}\otimes
p_2^*\mathcal L^{-1}$. On the diagonal its canonical height is
that of $[2]^*\mathcal L\otimes\mathcal L^{-2}\simeq
\mathcal L^2$, namely $2\widehat h_O$ in the absolute convention.
Changing from the absolute convention to relative heights over
$K$ multiplies by $[K:\Q]=2$. Therefore
\[
 \widehat h_K(Y)=4\widehat h_O(Y_d)
                =a_i^2\widehat h_O(P_i).
\]
This factor comes from the Poincar\'e pairing and the base field,
independently of any numerical height routine.

\item A real-adapted basis of $\Lambda_0$ gives
\[
 \Omega_i\Omega_{i,D}=
 \begin{cases}
 2V_0/(|d|\sqrt{|D|}),&\Delta(E_0)<0,\\
 4V_0/(|d|\sqrt{|D|}),&\Delta(E_0)>0.
 \end{cases}
\]
Indeed, a triangular real lattice has one real component and its
purely imaginary primitive period has twice the height of its
nonreal basis vector. A rectangular lattice has two real
components. Twisting by the two opposite signs interchanges the
real and imaginary directions, and the exact minimal differential
scalings in step 2 supply $1/(|d|\sqrt{|D|})$.
Substituting into \eqref{ari:eq:Cgz-product} gives
\begin{equation}\label{ari:eq:Cgz-denominator}
 q_ir_i=
 \begin{cases}
 a_i^2/c_0^2,&\Delta(E_0)<0,\\
 a_i^2/(2c_0^2),&\Delta(E_0)>0.
 \end{cases}
\end{equation}
In particular, if $r_i=A_i/B_i'$ is reduced and $c_0=1$, the
denominator of $q_i$ divides $2|A_i|$. We use the weaker upper
bound $B_i=4|A_i|$ in the interval certificate.

\item \paragraph{4. Global identification of the CM coefficient series.}
For $s=11,43,67,163$, put
$F_s=\Q(\sqrt{-s})$ and $\mathfrak p=(\sqrt{-s})$.
These fields have class number one and units $\{\pm1\}$.
The character
\[
 \psi((\alpha))=
  \left(\frac{\alpha\bmod\mathfrak p}{s}\right)\alpha
\]
is well defined because $(-1/s)=-1$. It has infinity type $(1,0)$
and finite conductor $\mathfrak p$. Hecke's theta construction
places $\sum_{\mathfrak a}\psi(\mathfrak a)q^{N\mathfrak a}$ in
$S_2(\Gamma_0(s^2))$: the nebentypus is
$\chi_{-s}(n)(n/s)=1$. This exact special case is also written
in \citet[\S3, pp.~4--5]{CrDarmon2009}. Grouping conjugate
generators $\alpha=(u+v\sqrt{-s})/2$ gives the integral formula
used by the coefficient stream:
\[
 b_n=\sum_{a^2=n,\ a>0}\left(\frac as\right)a
  +\sum_{\substack{u,v>0,\ u\equiv v\ (2)\\u^2+sv^2=4n}}
       \left(\frac{(s+1)u/2}{s}\right)u.
\]
For all four bases, exact PARI elliptic-curve coefficients equal
$b_n$ for $1\leq n\leq10000$.

\item This finite comparison proves equality of the entire modular forms.
For completeness, if a nonzero form $g$ of weight $k$ on
$\Gamma_0(N)$ has trivial character, take the product of
$g|_k\gamma$ over right coset representatives in
$\Gamma_0(N)\backslash\mathrm{SL}_2(\Z)$. It is a nonzero
holomorphic full-level form of weight $kI$, where
$I=N\prod_{p\mid N}(1+1/p)$. Every factor has nonnegative
order at infinity, so the product's order is at least
$\operatorname{ord}_\infty g$. The full-level valence formula
therefore implies $\operatorname{ord}_\infty g\leq kI/12$.
Equivalently, agreement through $\lfloor kI/12\rfloor$ proves
equality. The stronger arithmetic congruence statement is
\citet[Theorem 9.18]{CgzStein2007}; only the elementary
characteristic-zero argument is needed here. For $k=2,N=s^2$
the four bounds are
\[
 \left\lfloor\frac{s^2+s}{6}\right\rfloor
     =22,315,759,4455,
\]
all below $10000$. Applying the exact quadratic twist character
then identifies every coefficient in both the target and auxiliary
streams, including coefficients past the tested range. The exact
metadata transcript prints the four bounds and their comparisons.

\item \paragraph{5. Directed enclosures of both $L$-sums.}
For $c=2\pi/\sqrt N$, the root-number $+1$ formula is
\[
 L(E,1)=2\sum_{n\geq1}a_n\frac{e^{-cn}}n,
\]
while for root number $-1$ the first derivative is
\eqref{ari:eq:Cr-first}. The exact coefficients in the smaller cases
come from retained PARI files. The largest auxiliary stream uses
$s=163$, signed twist $-557403$, conductor
$8254937936042721$, and $M=867617471$; the target stream uses
$s=163,d=185801,M=289205824$. Both use the cutoff $cM\geq60$
up to rounding of the integer cutoff, with the actual $M$ used in
the tail bound. No precision conclusion depends on the nominal
cutoff.

\item Here are the finite-sum and rounding certificates implemented in
\path{code/B_C6_central_stream.cpp}. All intervals use 256-bit
MPFR with outward rounding at every operation. In a block centred
at $b$ with radius $R=128$, set
$S_k=\sum a_n(n-b)^k\in\Z$. For $g(z)=e^{-cz}/z$, put
\[
 C_0=1,\quad C_k=C_{k-1}c/k,\qquad
 T_0=1/b,\quad T_k=(T_{k-1}+C_k)/b.
\]
The coefficient of $(z-b)^k$ is $(-1)^ke^{-cb}T_k$.
For $h(z)=E_1(cz)/z$, put
\[
 e_0=E_1(cb),\quad e_k=(-1)^ke^{-cb}T_{k-1}/k\ (k\geq1),
 \qquad h_0=e_0/b,\quad h_k=(e_k-h_{k-1})/b.
\]
The block sums are the finite sums of these interval coefficients
times the exact integers $S_k$.

\item On $|z-b|=b/2$, the two kernels have bounds
\[
 |g(z)|\leq\frac{2e^{-cb/2}}b=:U_g,
 \qquad
 |h(z)|\leq\frac{4e^{-cb/2}}{cb^2}=:U_h.
\]
For the second bound, use
$E_1(cz)=\int_1^\infty e^{-czt}\,dt/t$ when $\Re z>0$.
Thus truncation at degree $k$ contributes at most
\[
 U_g\frac{(2R/b)^{k+1}}{1-2R/b}\sum|a_n|
 \quad\hbox{or}\quad
 U_h\frac{(2R/b)^{k+1}}{1-2R/b}\sum|a_n|.
\]
The degree is selected to make the geometric factor small, but
the actual remainder is separately evaluated with upward rounding.
For the derivative stream, the initial $E_1$ in each segment is
evaluated using correctly rounded $-\operatorname{Ei}(-x)$ and
monotonicity. Subsequent centres, spaced by $W=257$, use
\[
 E_1(c(b+W))=E_1(cb)-\int_0^Wg(b+x)\,dx.
\]
Integrating the same Taylor expansion gives the interval for this
advance. Its omitted part is bounded by
$WU_g(2W/b)^{k+1}/(1-2W/b)$ and added to both endpoints.
The first $10280$ terms are summed directly. Segment boundaries
are multiples of $257$, so no coefficient is omitted or repeated.

\item Hasse's bound and multiplicativity give $|a_n|\leq2n$.
The infinite tails, including the factor two, are consequently at
most
\[
 \frac{4e^{-c(M+1)}}{1-e^{-c}},\qquad
 \frac{4e^{-c(M+1)}}{c(M+1)(1-e^{-c})},
\]
respectively. The program evaluates both with the downward bound
for $c$ and upward bounds for the remaining positive expressions.
It finally adds the block remainders, interval arithmetic errors,
and infinite tail. For row eleven the central-value tail is below
$5.07\cdot10^{-19}$, and the derivative tail is below
$2.82\cdot10^{-21}$. These modest precisions already suffice for
rational separation. All twelve completed interval outputs are
retained, as are the exact coefficient files.

\item \paragraph{6. Periods, heights, and exact rational separation.}
The periods are enclosed independently with 512-bit Arb. Let
$\tau=i\sqrt s$ for $i=2,4$ and
$\tau=(1+i\sqrt s)/2$ for $i=6,9,10,11$; here
$s=3,7,11,43,67,163$, respectively. The exact CM identification
and reduced-form class-number check are part of \ref{ari:C.6.5}. Put
$T=\Im\tau$, $q_\tau=e^{2\pi i\tau}$, and
$r=|q_\tau|<10^{-4}$. We enclose
\[
 E_4(\tau)=1+240\sum_{n=1}^{100}\sigma_3(n)q_\tau^n
                   +\text{tail}.
\]
Since $\sigma_3(n)<2n^3$, with $A=101$ the absolute tail is at most
\[
 480r^A\left\{
 \frac{A^3}{1-r}+\frac{3A^2r}{(1-r)^2}
 +\frac{3Ar(1+r)}{(1-r)^3}
 +\frac{r(1+4r+r^2)}{(1-r)^4}\right\}.
\]
The exact invariants $c_4,c_6$ of each minimal equation give
$\kappa=2\pi(E_4(\tau)/c_4)^{1/4}$. The sign of $E_6(\tau)$
is positive: $\sigma_5(n)<2n^5$ and
$n^5\leq32^{n-1}$ give
$|E_6(\tau)-1|\leq1008r/(1-32r)<1$.
Thus the full real period is
\[
 \Omega=n_{\R}\kappa\begin{cases}
 1,&c_6>0,\\
 T,&c_6<0,\ \Delta>0,\\
 2T,&c_6<0,\ \Delta<0,
 \end{cases}
 \qquad n_{\R}=\begin{cases}2,&\Delta>0,\\1,&\Delta<0.\end{cases}
\]
These formulas apply separately to the target and auxiliary curve.
Every operation, the Fourier tail, and the fourth root is enclosed
by Arb. The height intervals are the independent certificates
\path{records/C_cm_height_arb.json} of \ref{ari:C.6.5}.

\item The final computation is
\path{code/B_C6_exact_factors.py}. First, the intervals for
$r_i$ isolate a unique integer after multiplication by the exact
$t_bn_b$ in \eqref{ari:eq:Cgz-central-integrality}. This proves the
six auxiliary rational values, with no BSD input. Next form the
interval $L'(E_i,1)/(\Omega_i\widehat h_O(P_i))$. The enclosures
below are deliberately enlarged decimal summaries of the recorded
Arb intervals:
\[
\begin{array}{c|c|c}
 i&|r_i-r_i^{\rm listed}|\text{ is at most}
   &|q_i-q_i^{\rm listed}|\text{ is at most}\\\hline
 2&4\cdot10^{-22}&2\cdot10^{-53}\\
 4&4\cdot10^{-20}&3\cdot10^{-45}\\
 6&4\cdot10^{-22}&2\cdot10^{-53}\\
 9&5\cdot10^{-21}&4\cdot10^{-53}\\
 10&3\cdot10^{-20}&10^{-23}\\
 11&7\cdot10^{-15}&8\cdot10^{-19}
\end{array}
\]
Each last-column bound is strictly less than $1/B_i$. A rational
number with denominator at most $B_i$ differing from the listed
integer would differ by at least $1/B_i$. Equation
\eqref{ari:eq:Cgz-denominator} therefore makes every listed $q_i$
exact. The same equation gives the listed $|a_i|$ by integer
square-root checks. The JSON certificate records all input
intervals, exact invariants, rational values, and positive separation
margins. Its stdout concludes that all six remaining factors are
exactly certified. This completes the proof.
\end{enumerate}
\end{proof}


\subsection{Supplement: the three specified smooth torsion fibres}

\begin{proposition}[T1; Proved, with exact local and descent computations]\label{ari:C.6.T1}
\StatementRecord{ari:C.6.T1}{proved}{Historical events: GPT-6 Astra (OpenAI), ChatGPT Pro mode, effort Pro, S6\_09\_PAPERS, 17 November 2026: targeted proof check. Also checked in this writing. Pre-continuation writing events: GPT-6 Astra (OpenAI), ChatGPT Pro mode, effort Pro, S6\_10\_PAPERS, 18 November 2026: computation or formal certificate; targeted proof check. Continuation metadata events: GPT-6 Astra (OpenAI), ChatGPT Pro mode, effort Pro, S6\_10\_PAPERS continuation, 19 November 2026: used and rechecked in this run; status-vocabulary reconciliation; metadata reconciliation only; no new mathematical check. \emph{Pre-continuation recorded proof scope.} Read complete printed arithmetic proofs and re-derived the exact scaling/local-height equations and the genus-trace square denominator argument. Rank, saturation, directed L-derivative replays and precise external theorem matches remain explicit gates; no proof verdict is awarded to those computational claims by this read. Read and replayed the unmodified local and final torsion programs in PARI2.17.2. Exact rank0,torsion4,conductor40,discriminant80,component orders2/1 and local-height coefficients1/12 match. Exact plus modular symbol is-1/4; the printed period/L-value comparisons reproduce. All starred final-log entries are stack-growth warnings; complete transcript contains no GP error.}{the hypotheses, notation and definitions stated here; the explicit constructions and calculations in the proof}

At $t=32/5$ a global minimal model and the image of the section are
\[
 E:y^2=x^3-2x+1,\qquad P=(0,1),\qquad \Delta=80.
\]
The conductor is $40$, $E(\Q)=\Z/4\Z$ is generated by $P$,
and $\widehat h(P)=0$. The index of $\Z P$ in the quotient
modulo torsion is one. The complete bad-prime data are
\[
\begin{array}{c|c|c|c|c}
p&\operatorname{ord}_p\Delta&\text{type}&\text{component of }P&\lambda_p(P)\\\hline
2&4&\mathrm{III}&\text{unique nonzero class, order }2&\log2/12\\
5&1&\mathrm I_1&\text{identity}&\log5/12.
\end{array}
\]
All other finite local heights vanish. Consequently
\[
 \log N=\frac{\log10}{12},\qquad
 D_{\rm part}=\frac6\pi\log\frac{10^{1/12}}{2\sqrt3}.
\]
Numerically the two sides of the global height identity equal
\[
 0.191882091082837140334832621223697017300091790719064414669444\ldots.
\]
With the full real period $\Omega=5.9376498936895458116293755517\ldots$,
the central derivative is
\[
 \frac{L'(E,1)}\Omega=
 0.071331625531738774092617964261659879861112925567606085287841\ldots.
\]
The analytic rank is zero; a vanishing first derivative is therefore
not expected even though the point $P$ is torsion.
\end{proposition}


\begin{proof}
\begin{enumerate}[label=\arabic*.,leftmargin=*]
\item The change of variables from the displayed family is
$x=(24/5)^2X$, $y=(24/5)^3Y$. Tate's algorithm gives the
listed minimal model, conductor and local data. Since $P$ has
order four, $3P=-P$. The local division identity consequently
gives
\[
 \lambda_p(P)=\left(\frac{v_p(\Delta)}{12}
                         -\frac{v_p(\psi_3(P))}{8}\right)\log p.
\]
Here $\psi_3(P)=-4$, proving both local coefficients exactly.
The exact descent output is $[0,0,0,[]]$, and the torsion
calculation gives $\Z/4\Z$. These computations, the direct
global equality and the analytic derivative are in
\path{code/Crank_torsion_local.gp} and
\path{code/Crank_torsion_final.gp}, with their outputs.
As a further check, the exact plus modular symbol from $0$ to
$\infty$ is $-1/4$ in the real-period lattice normalization;
the numerical ratio $L(E,1)/\Omega$ is $1/8$. The completed
functional equation then gives the checked expression
$(\gamma+\log(2\pi/\sqrt{40}))/8$ for the displayed derivative
ratio. The local-height proof does not use this $L$-value check.
\end{enumerate}
\end{proof}

\begin{proposition}[T2; Proved, with exact number-field local arithmetic]\label{ari:C.6.T2}
\StatementRecord{ari:C.6.T2}{proved}{Historical events: GPT-6 Astra (OpenAI), ChatGPT Pro mode, effort Pro, S6\_09\_PAPERS, 17 November 2026: targeted proof check. Also checked in this writing. Pre-continuation writing events: GPT-6 Astra (OpenAI), ChatGPT Pro mode, effort Pro, S6\_10\_PAPERS, 18 November 2026: computation or formal certificate; targeted proof check. Continuation metadata events: GPT-6 Astra (OpenAI), ChatGPT Pro mode, effort Pro, S6\_10\_PAPERS continuation, 19 November 2026: used and rechecked in this run; status-vocabulary reconciliation; metadata reconciliation only; no new mathematical check. \emph{Pre-continuation recorded proof scope.} Read complete printed arithmetic proofs and re-derived the exact scaling/local-height equations and the genus-trace square denominator argument. Rank, saturation, directed L-derivative replays and precise external theorem matches remain explicit gates; no proof verdict is awarded to those computational claims by this read. Read and replayed the exact local program and the Simon two-isogeny descent driver with only its library read paths changed. bnfcertify=1 and LIMBIGPRIME=0; the complete new transcript has the displayed surviving points and excluded local coverings, quotient orders2 and4, hence rank1. Torsion6,three local types,valuations and weighted height sum match. Also reproduced the optional100-digit number-field L-value display, which retains numerical scope. No independent review of the entire Simon library is claimed.}{the hypotheses, notation and definitions stated here; \cref{ari:C.4.4}}

Put $s=\sqrt3$ and
$t_\pm=(3008\pm1728s)/11$. Over $K=\Q(s)$, a global minimal
model for the fibre at $t_-$ is
\[
 E_-:y^2=x^3-(18+10s)x+21+12s,\qquad P_-=(0,3+2s).
\]
The model and point for $t_+$ are its conjugates. The
discriminant of this model is $1032912+596352s$. Define
\[
 \mathfrak p_2=(2,1+s),\qquad\mathfrak p_3=(s),\qquad
 \mathfrak p_{11,-}=(11,s+5),\quad
 \mathfrak p_{11,+}=(11,s-5).
\]
The bad-prime data of $E_-$ are
\[
\begin{array}{c|c|c|c|c|c}
v&e_v&v(\Delta)&\text{type}&\text{component order of }P_-&\lambda_v(P_-)\\\hline
\mathfrak p_2&2&8&\mathrm{IV}^*&3&0\\
\mathfrak p_3&2&3&\mathrm{III}&2&0\\
\mathfrak p_{11,-}&1&1&\mathrm I_1&1&\log11/12.
\end{array}
\]
For $E_+$ replace $\mathfrak p_{11,-}$ by
$\mathfrak p_{11,+}$. The nonzero component at
$\mathfrak p_3$ is unique; the two nonzero classes at
$\mathfrak p_2$ are interchanged by $P\mapsto-P$ and the
class of $P$ is a generator. All omitted local heights vanish.
Both points have order six, and
\begin{align*}
 \log N(P_-)&=\frac1{12}\log(17+10\sqrt3),\\
 \log N(P_+)&=\frac1{12}\log(10\sqrt3-17),\\
 \sum_{\sigma:K\hookrightarrow\R}\log N(P_\sigma)
 &=\frac{\log11}{12}
 =\sum_{v<\infty}n_v\lambda_v(P)-2\widehat h(P).
\end{align*}
The common value in the last line is
\[
 0.199824606066530878671828631497094108318475571161451431268214\ldots.
\]
In particular $A_\mp=(17+10\sqrt3)^{1/12}$ and
$(10\sqrt3-17)^{1/12}$ respectively give
$D_{\rm part}(t_\mp)=(6/\pi)\log(A_\mp/(2\sqrt3))$.

The natural Mordell--Weil group here is $E_\pm(K)$; the
notation $E_{t_\pm}(\Q)$ requested in the stated conventions does not
define an elliptic curve over $\Q$, since its $j$-invariant
is irrational. Over $K$, both curves have rank one and torsion
group $\Z/6\Z$. Thus the index of the torsion section in the
free quotient is infinite. For example $(1,-1-s)$ has infinite
order on $E_-$.
\end{proposition}


\begin{proof}
\begin{enumerate}[label=\arabic*.,leftmargin=*]
\item For $t_-$ the change from the original family to this model
has scale $u=48(-64+37s)/11$, with translation parameters zero.
It is obtained by multiplying the first global minimalization
scale by the unit $(2+s)^{-3}$. Hence it preserves all finite
minimal valuations. The conjugate scale gives $E_+$.
Tate's algorithm yields the displayed table. Its discriminant
ideal is
$\mathfrak p_2^8\mathfrak p_3^3\mathfrak p_{11,-}$ and its
conductor ideal is
$\mathfrak p_2^2\mathfrak p_3^2\mathfrak p_{11,-}$.

\item As $5P=-P$, the normalized local heights are computed exactly by
\[
 \frac{\lambda_v(P)}{\log p}
 =\frac1{e_v}\left(\frac{v(\Delta)}{12}
                         -\frac{v(\psi_5(P))}{24}\right).
\]
The valuations of $\psi_5(P)$ at the three listed primes are
$16,6,0$, respectively. This proves the table, including zero
contribution at the two additive primes. Here $f_v=1$ at all
three primes, so weighting by $n_v=e_vf_v$ gives precisely
$\log11/12$. The exact torsion values proved in Proposition~\ref{ari:C.4.4} have product
$(17+10\sqrt3)(10\sqrt3-17)=11$, proving the archimedean side.

\item For the rank, translate the rational two-torsion point to zero.
The resulting equation is
\[
 y^2=x^3-(9+6s)x^2+(45+26s)x.
\]
The $2$-isogeny descent gives
$[E(K):\widehat\phi E'(K)]=2$ and
$[E'(K):\phi E(K)]=4$. Every surviving covering has a
displayed point; every excluded covering has a recorded local
obstruction at $\mathfrak p_2$ or $\mathfrak p_3$.
The quotient formula $2^{r+2}=2\cdot4$ proves $r=1$.
The field arithmetic was certified by \texttt{bnfcertify};
Simon descent was run with \texttt{LIMBIGPRIME=0}, so all local
tests were deterministic. The complete transcript is
\path{code/Crank_torsion_descent_sage.out}. The unmodified
Simon scripts maintained in Sage are retained with their source
URLs; the older upstream scripts and their unsuccessful
compatibility attempts are also retained. The torsion calculation
gives $\Z/6\Z$, generated by $P$, independently of the rank
calculation.
\end{enumerate}
\end{proof}

\paragraph{The corresponding number-field $L$-derivative.}
For completeness, the natural replacement for the nonexistent
rational-curve quantity at $t_\pm$ can also be computed. Put
\[
 \Omega_K=\frac1{\sqrt{12}}
  \prod_{\sigma:K\hookrightarrow\R}
       \int_{E_\sigma(\R)}|\omega_\sigma|.
\]
The Hasse--Weil $L$-function over $K$ has analytic conductor
$12^2\cdot396=57024$, gamma parameters $(0,0,1,1)$, and sign
$-1$. Direct 100-digit evaluation gives
\begin{align*}
 L'(E/K,1)&=2.046708855992864213402012895411307398829979893338038716037006\ldots,\\
 L'(E/K,1)/\Omega_K
 &=0.260797680551692455136263550519918387218960539765929916318910\ldots.
\end{align*}
This is a value for the curve over $K$, common to the two
conjugates. If the period is instead defined as the product of
the two full real periods, omitting the discriminant factor,
the ratio is
$0.075285805535274833467509818428226475613981023404668273560620\ldots$.
Both conventions and all input periods are explicitly recorded
in \path{code/Crank_torsion_final.out}; neither value enters
the exact torsion height proof.



\subsection{Integral and Gaussian descriptions}

\begin{proposition}[the Mahler-measure identities]\label{ari:C.8.6}
\StatementRecord{ari:C.8.6}{proved}{Historical events: Claude Fable 5.2 (Anthropic), effort high, 16 November 2026: used and rechecked in this run. GPT-6 Astra (OpenAI), ChatGPT Pro mode, effort Pro, S6\_09\_PAPERS, 17 November 2026: targeted proof check. Also checked in this writing. Pre-continuation writing events: GPT-6 Astra (OpenAI), ChatGPT Pro mode, effort Pro, S6\_10\_PAPERS, 18 November 2026: targeted proof check. Continuation metadata events: GPT-6 Astra (OpenAI), ChatGPT Pro mode, effort Pro, S6\_10\_PAPERS continuation, 19 November 2026: used and rechecked in this run; status-vocabulary reconciliation; metadata reconciliation only; no new mathematical check. \emph{Pre-continuation recorded proof scope.} Read complete proof and re-derived its Green-function divisor/mean, beta integral, signed Gaussian Fourier, Betti Jacobian/current-density, division-divisor multiplicity, or heat-equation calculation. For C.8.10 the DeMarco-Mavraki external theorem locator remains a separate gate.}{the hypotheses, notation and definitions stated here; \cref{ari:C.8.5}}

Let $\Lambda\subset\C$ be a lattice of area $A$, let $v$ be a zero
of $\wp_\Lambda$, and let $d\mu=A^{-1}dA_z$ be Haar probability
measure on $\C/\Lambda$. Put $U(z)=\log N(z)$ and
$L=U(v)$. Then
\begin{align}
 6L&=\int_{\C/\Lambda}\log\left|\frac{\wp'(z)^2}{\wp(z)^3}\right|d\mu(z),
 \label{ari:eq:ca-mahler}\\
 N(v)^2&=|\Delta(\Lambda)|^{1/6}
      \exp\left(-\int_{\C/\Lambda}\log|\wp(z)|d\mu(z)\right).
 \label{ari:eq:ca-wpmean}
\end{align}
If $C(u)=1-(g_2/4)u^2-(g_3/4)u^3$, the pushforward of $d\mu$
under $u=1/\wp(z)$ is
\begin{equation}\label{ari:eq:ca-nu}
 d\nu(u)=\frac1{2A}\frac{dA_u}{|u|\,|C(u)|}.
\end{equation}
In particular its unnormalized integral is exactly $2A$, and
\[
 N(v)^6=4\exp\left(\int_{\C}\log|C(u)|d\nu(u)\right).
\]
Writing $e_1,e_2,e_3$ for the three branch values and
$p(x)=\int\log|x-\wp(z)|d\mu(z)$, the maximum theorem is equivalently
\[
 \int\log|C|d\nu\leq\log2,\qquad
 p(0)\geq\frac{p(e_1)+p(e_2)+p(e_3)}3-\frac{\log2}3.
\]
Equality in either inequality occurs exactly when $g_3=0$.
\end{proposition}


\begin{proof}
\begin{enumerate}[label=\arabic*.,leftmargin=*]
\item The mean of $U$ is zero by the theta product calculation in C.1,
and near zero
$U(z)=\log|z|+\frac1{12}\log|\Delta|+o(1)$.
The divisor of $\wp$ is $(v)+(-v)-2(O)$, with multiplicity
two when $v=-v$. Subtract its logarithmic divisor from the
corresponding combination of $U$'s. The result is harmonic
on the compact torus and hence constant. Its expansion at
zero gives
\[
 \log|\wp(z)|=U(z-v)+U(z+v)-2U(z)
                  +\frac16\log|\Delta|-2L.
\]
Integration proves \eqref{ari:eq:ca-wpmean}.
Similarly the divisor of $\wp'$ is the sum of the three
nonzero half-periods minus $3(O)$. The multiplication-by-two
identity of C.1 gives the product of $N$ at those half-periods
as $2$. Since $\wp'(z)\sim-2z^{-3}$, comparison at zero yields
\[
 \log|\wp'(z)|=\sum_{w\in E[2]\setminus\{O\}}U(z-w)-3U(z)
                         +\frac14\log|\Delta|.
\]
Its mean is $\frac14\log|\Delta|$. Twice this identity minus
three times the preceding mean proves \eqref{ari:eq:ca-mahler}.
All logarithmic singularities are locally integrable, so these
integrations also cover a zero at a half-period.

\item The map $z\mapsto x=\wp(z)$ has degree two. Its two inverse
sheets each have Jacobian $|\wp'(z)|^{-2}$; branch values form
a measure-zero set. Thus its pushforward density is
$2/(A|4x^3-g_2x-g_3|)$. Under $x=1/u$ the Jacobian is
$|u|^{-4}$ and the denominator is $4|u|^{-3}|C(u)|$.
This proves \eqref{ari:eq:ca-nu} and its normalization. The identity
$\wp'^2/\wp^3=4C(u)$ gives the next formula.

\item Replacing $v$ above by a half-period gives
$p(e_j)=\frac16\log|\Delta|-2U(w_j)$; summing and using
$\sum_jU(w_j)=\log2$ proves the centroid formulation.
The equality characterization is that of \ref{ari:C.8.5},
since $j=1728$ is equivalent to $g_3=0$ for a nonsingular lattice.
\end{enumerate}
\end{proof}

\begin{proposition}[the two complex beta integrals]\label{ari:C.8.7}
\StatementRecord{ari:C.8.7}{proved}{Historical events: Claude Fable 5.2 (Anthropic), effort high, 16 November 2026: used and rechecked in this run. GPT-6 Astra (OpenAI), ChatGPT Pro mode, effort Pro, S6\_09\_PAPERS, 17 November 2026: targeted proof check. Also checked in this writing. Pre-continuation writing events: GPT-6 Astra (OpenAI), ChatGPT Pro mode, effort Pro, S6\_10\_PAPERS, 18 November 2026: targeted proof check. Continuation metadata events: GPT-6 Astra (OpenAI), ChatGPT Pro mode, effort Pro, S6\_10\_PAPERS continuation, 19 November 2026: used and rechecked in this run; status-vocabulary reconciliation; metadata reconciliation only; no new mathematical check. \emph{Pre-continuation recorded proof scope.} Read complete proof and re-derived its Green-function divisor/mean, beta integral, signed Gaussian Fourier, Betti Jacobian/current-density, division-divisor multiplicity, or heat-equation calculation. For C.8.10 the DeMarco-Mavraki external theorem locator remains a separate gate.}{the hypotheses, notation and definitions stated here; \cref{ari:C.8.6}; \cref{ari:eq:ca-mahler}}

For the square and equianharmonic lattices respectively, the
integrals in \ref{ari:C.8.6} are
\[
 \int\log|C|\,d\nu=\log2,\qquad
 \int\log|C|\,d\nu=\frac32\log3-2\log2.
\]
They recover $N^2=2$ and $N^4=3$ with the same normalization.
\end{proposition}


\begin{proof}
\begin{enumerate}[label=\arabic*.,leftmargin=*]
\item By homogeneity it suffices to take $C(u)=1-u^m$, with $m=2$
or $m=3$. For $-1<\Re s<1-1/m$ put
\[
 I_m(s)=\int_{\C} |u|^{-1}|1-u^m|^{s-1}dA_u.
\]
The substitution $w=u^m$, with all $m$ sheets counted, gives
\[
 I_m(s)=\frac1m B_{\C}\left(\frac1{2m},\frac{s+1}2\right),
 \quad
 B_{\C}(a,b)=\pi\frac{\Gamma(a)\Gamma(b)\Gamma(1-a-b)}
                   {\Gamma(1-a)\Gamma(1-b)\Gamma(a+b)}.
\]
We include a proof of the integral evaluation. For
$0<a,b$ and $a+b<1$, insert
\[
 |w|^{2a-2}=\frac1{\Gamma(1-a)}
       \int_0^\infty t^{-a}e^{-t|w|^2}\,dt
\]
and the corresponding formula for $|1-w|^{2b-2}$.
Tonelli's theorem applies to the positive integrand. The
Gaussian integral over $w$ is
$\pi(t+s)^{-1}\exp(-ts/(t+s))$.
Putting $t=rv$, $s=r(1-v)$ reduces the remaining integrals
to $\Gamma(1-a-b)$ and the beta integral $B(a,b)$, proving
the displayed formula. The same integrable bounds justify
differentiation in a neighbourhood of $s=0$.

\item Consequently, with $a=1/(2m)$ and $\psi=\Gamma'/\Gamma$,
\[
 \frac{I_m'(0)}{I_m(0)}
 =\frac12\left(2\psi(1/2)-\psi(1/2-a)-\psi(1/2+a)\right).
\]
The gamma duplication and triplication formulas, differentiated
logarithmically, evaluate this as $\log2$ for $m=2$ and
$\frac32\log3-2\log2$ for $m=3$. These can also be checked
directly from
$\psi(1/2)=-\gamma-2\log2$,
$\psi(1/4)+\psi(3/4)=-2\gamma-6\log2$ and
$\psi(1/3)+\psi(2/3)=-2\gamma-3\log3$.
Equation \eqref{ari:eq:ca-mahler} completes the normalization test.
\end{enumerate}
\end{proof}

\begin{proposition}[the signed Gaussian identity]\label{ari:C.8.8}
\StatementRecord{ari:C.8.8}{proved}{Historical events: Claude Fable 5.2 (Anthropic), effort high, 16 November 2026: used and rechecked in this run. GPT-6 Astra (OpenAI), ChatGPT Pro mode, effort Pro, S6\_09\_PAPERS, 17 November 2026: targeted proof check. Also checked in this writing. Pre-continuation writing events: GPT-6 Astra (OpenAI), ChatGPT Pro mode, effort Pro, S6\_10\_PAPERS, 18 November 2026: targeted proof check. Continuation metadata events: GPT-6 Astra (OpenAI), ChatGPT Pro mode, effort Pro, S6\_10\_PAPERS continuation, 19 November 2026: used and rechecked in this run; status-vocabulary reconciliation; metadata reconciliation only; no new mathematical check. \emph{Pre-continuation recorded proof scope.} Read complete proof and re-derived its Green-function divisor/mean, beta integral, signed Gaussian Fourier, Betti Jacobian/current-density, division-divisor multiplicity, or heat-equation calculation. For C.8.10 the DeMarco-Mavraki external theorem locator remains a separate gate.}{the hypotheses, notation and definitions stated here; the explicit constructions and calculations in the proof}

For $\Lambda=\Z+\tau\Z$, $T=\Im\tau$, and every $z\in\C$,
\[
 \sqrt{2T}\,|\eta(\tau)|^2N(z,\tau)^2
 =\sum_{m\in\frac12\Lambda}\epsilon(m)
                         e^{-2\pi|z-m|^2/T},
\]
where $\epsilon=-1$ on $\Lambda$ and $+1$ on the three other
cosets of $\Lambda$ in $\frac12\Lambda$.
Both sides and all their derivatives converge locally uniformly.
\end{proposition}


\begin{proof}
\begin{enumerate}[label=\arabic*.,leftmargin=*]
\item Write $z=b+a\tau$ and expand in the Fourier characters
$e^{2\pi i(kb+la)}$ of $(a,b)\in(\R/\Z)^2$.
For $G(z)=\sum_{\lambda\in\Lambda}e^{-2\pi|z-\lambda|^2/T}$,
Poisson summation gives the coefficient
\[
 \frac12\exp\left(-\frac{\pi|k\tau-l|^2}{2T}\right).
\]
The right side is the sum of the four half-period translates
of $G$, minus $2G$. Its coefficient is therefore
\[
 (-1)^{k+l+kl}
       \exp\left(-\frac{\pi|k\tau-l|^2}{2T}\right).
\]
To check the sign, it is $+1$ when both $k,l$ are even and
$-1$ otherwise.

\item On the left use
\[
 \vartheta_1(z,\tau)
 =-i\sum_{n\in\Z}(-1)^n
        e^{\pi i(n+1/2)^2\tau+2\pi i(n+1/2)z}.
\]
Multiply by its conjugate and by
$\sqrt{2T}e^{-2\pi(\Im z)^2/T}$.
In the double sum put $k=n-m$ and apply one-dimensional
Poisson summation to the remaining Gaussian in $m$.
Completing the square gives precisely the coefficient just
displayed. This proves equality of the Fourier series.
Positive definiteness of $|k\tau-l|^2$ and Gaussian decay
justify the changes of order, Poisson summation, and every
claimed differentiation.
\end{enumerate}
\end{proof}




\subsection{The Betti measure and torsion equidistribution}

We take $dd^c=(i/\pi)\partial\bar\partial$, so that
$dd^c\log|w|=\delta_0$. For a local complex coordinate $w$,
$dA_w=(i/2)dw\wedge d\bar w$.

\begin{proposition}[Refuted sign and normalization, with the corrected current identity proved]\label{ari:C.8.9}
\StatementRecord{ari:C.8.9}{proved}{Historical events: Claude Fable 5.2 (Anthropic), effort high, 15 November 2026: carried as background without fault found. Claude Fable 5.2 (Anthropic), effort high, 16 November 2026: used and rechecked in this run. GPT-6 Astra (OpenAI), ChatGPT Pro mode, effort Pro, S6\_09\_PAPERS, 17 November 2026: targeted proof check. Also checked in this writing. Pre-continuation writing events: GPT-6 Astra (OpenAI), ChatGPT Pro mode, effort Pro, S6\_10\_PAPERS, 18 November 2026: targeted proof check. Continuation metadata events: GPT-6 Astra (OpenAI), ChatGPT Pro mode, effort Pro, S6\_10\_PAPERS continuation, 19 November 2026: used and rechecked in this run; status-vocabulary reconciliation; metadata reconciliation only; no new mathematical check. \emph{Pre-continuation recorded proof scope.} Read complete proof and re-derived its Green-function divisor/mean, beta integral, signed Gaussian Fourier, Betti Jacobian/current-density, division-divisor multiplicity, or heat-equation calculation. For C.8.10 the DeMarco-Mavraki external theorem locator remains a separate gate.}{the hypotheses, notation and definitions stated here; \cref{ari:C.3.2}; \cref{ari:eq:ca-cuspexp}}

On a local period chart write $\mu=a\tau+b$, with $a,b$ real,
and put $T=\Im\tau$. The positive Betti measure is
\[
 m_P=db\wedge da
 =\frac{|\mu_w-a\tau_w|^2}{T}\,dA_w.
\]
Thus the proposed form $da\wedge db$ has the wrong sign.
The globally single-valued function $L(t)=\log N(P_t)$ satisfies
on $\mathbb P^1(\C)$
\begin{equation}\label{ari:eq:ca-current}
 dd^cL=\frac1{12}\delta_\infty-m_P,\qquad
               \int_{\mathbb P^1}m_P=\frac1{12}.
\end{equation}
There are no atoms of $m_P$ at $0$, $1$ or $\infty$.
Let $T_0=\sqrt3/2$, $\rho=(1+i\sqrt3)/2$, and
$p=\wp'_\rho((2-\rho)/3)$. The density in the ordinary
$t$ coordinate satisfies
\begin{align}
 \frac{dm_P}{dA_t}
 &=\frac{|t|^{-2/3}}{64T_0|p|^{2/3}}
                      \bigl(1+O(|t|^{1/3})\bigr),
                         &&t\longrightarrow0,\label{ari:eq:ca-density0}\\
 \frac{dm_P}{dA_t}
 &=\frac{|t-1|^{-3/2}}{24\pi^2\sqrt{E_4(i)}}
                      \bigl(1+O(|t-1|^{1/2})\bigr),
                         &&t\longrightarrow1.\label{ari:eq:ca-density1}
\end{align}
\end{proposition}


\begin{proof}
\begin{enumerate}[label=\arabic*.,leftmargin=*]
\item On a chart avoiding a zero of the section, the theta quotient
is holomorphic and nonzero. Its log modulus is pluriharmonic.
Only $-\pi(\Im\mu)^2/T$ contributes to the curvature.
Direct differentiation of this real quotient gives
\[
 \partial_w\partial_{\bar w}\frac{(\Im\mu)^2}{T}
       =\frac{|\mu_w-a\tau_w|^2}{2T}.
\]
Also
\[
 a_w=\frac{\mu_w-a\tau_w}{2iT},\qquad
 b_w=-\bar\tau\frac{\mu_w-a\tau_w}{2iT}.
\]
Taking the wedge product proves the formula for $m_P$ and
$dd^cL=-m_P$ on the chart. The transformation of $(a,b)$
under an integral change of period basis has determinant one,
so the measure patches. Translation of $\mu$ by a locally
constant integral period changes neither the formula nor
the measure. The formula remains valid on a root chart
at an elliptic modulus.

\item At $t=0$ the expansion proved in Proposition~\ref{ari:C.3.2} is
\[
 L(t)=\frac14\log3+\frac7{192}\Re t
       -B|t|^{4/3}+O(|t|^{5/3}),\qquad
 B=\frac{9\pi}{512T_0|p|^{2/3}}.
\]
The expansion is a convergent real-analytic expansion on a
root chart, so it can be differentiated termwise with the
indicated orders. The Euclidean identity
$\Delta_{\R^2}|t|^\alpha=\alpha^2|t|^{\alpha-2}$ gives
\eqref{ari:eq:ca-density0}.

\item At $t=1$ put $h=\mu-(1-\tau)/2$. The expansion in Proposition~\ref{ari:C.3.1}
and the modular $j$ expansion give
\[
 t-1=A_1h^4(1+O(h^2)),\qquad
 A_1=\frac94\pi^4E_4(i)>0,\qquad
 L=\log\sqrt2-\frac\pi2|h|^2+O(|h|^4).
\]
Therefore
$L=\log\sqrt2-\pi|t-1|^{1/2}/(2\sqrt{A_1})+O(|t-1|)$.
Termwise differentiation on the root chart proves
\eqref{ari:eq:ca-density1}. Both densities are locally integrable,
and the continuous potentials at $0$ and $1$ have no atoms.

\item For infinity take $u=1/t$, and set $\ell=\log(5+2\sqrt6)$.
The cusp branch is
$\mu=1/2-i\ell/(2\pi)+O(q)$, $q=e^{2\pi i\tau}$.
The theta and eta products, with the nonzero factor
$|2\sin\pi\mu|=2\sqrt3+O(q)$, give
\begin{equation}\label{ari:eq:ca-cuspexp}
 L=\frac1{12}\log|u|+\log(2\sqrt3)-\frac1{12}\log1728
          -\frac{\ell^2}{2\log(1728/|u|)}
          +O\left(\frac{|u|}{|\log|u||}+|u|\right).
\end{equation}
Here $q=u/1728+O(u^2)$; the estimate follows directly by
bounding the convergent products. The nonlogarithmic
correction is continuous at $u=0$ and has no point mass.
Equivalently its Betti density is
\[
 \frac{dm_P}{dA_u}
  =\frac{\ell^2}{2\pi|u|^2\log(1728/|u|)^3}
       \bigl(1+o(1)\bigr),
\]
which is integrable. The single logarithmic term in
\eqref{ari:eq:ca-cuspexp} contributes $\delta_\infty/12$.
This proves the global current identity; pairing it with
the constant function one gives the total mass.
\end{enumerate}
\end{proof}

\begin{proposition}[the correctly normalized equidistribution]\label{ari:C.8.10}
\StatementRecord{ari:C.8.10}{proved}{Historical events: Claude Fable 5.2 (Anthropic), effort high, 15 November 2026: carried as background without fault found. Claude Fable 5.2 (Anthropic), effort high, 16 November 2026: used and rechecked in this run. GPT-6 Astra (OpenAI), ChatGPT Pro mode, effort Pro, S6\_09\_PAPERS, 17 November 2026: targeted proof check. Also checked in this writing. Pre-continuation writing events: GPT-6 Astra (OpenAI), ChatGPT Pro mode, effort Pro, S6\_10\_PAPERS, 18 November 2026: targeted proof check; used and rechecked in this run. Continuation metadata events: GPT-6 Astra (OpenAI), ChatGPT Pro mode, effort Pro, S6\_10\_PAPERS continuation, 19 November 2026: used and rechecked in this run; status-vocabulary reconciliation; metadata reconciliation only; no new mathematical check. \emph{Pre-continuation recorded proof scope.} Read complete proof and re-derived its Green-function divisor/mean, beta integral, signed Gaussian Fourier, Betti Jacobian/current-density, division-divisor multiplicity, or heat-equation calculation. For C.8.10 the DeMarco-Mavraki external theorem locator remains a separate gate. Opened DeMarco–Mavraki Corollary1.2 p2, current pullback description pp13–14 and final equidistribution proof p23; rechecked the receiving passage from distinct Galois orbits to weighted division divisors using bounded fixed-point multiplicities. The probability limit is curvature divided by degree1/12, namely12m\_P; this closes the named source gate.}{the hypotheses, notation and definitions stated here; \citet[Theorem 1.1 and Corollary 1.2]{demarco2017}; \cref{ari:C.6.4}; \cref{ari:C.8.9}}

Let $D_n$ be the divisor, with intersection multiplicity, of
parameters $t\in\mathbb P^1\setminus\{0,1,\infty\}$ for which
$nP_t=O$. Then
\[
 \deg D_n=\frac{n^2}{12}+O(1),\qquad
 \frac1{n^2}[D_n]\ \longrightarrow\ m_P,\qquad
 \frac1{\deg D_n}[D_n]\ \longrightarrow\ 12m_P
\]
weakly as $n\to\infty$ through indices for which $D_n\ne0$.
For any nonrepeating sequence of individual algebraic torsion
parameters, their Galois-orbit probability measures likewise
converge to $12m_P$. Thus the probability limit requires the factor $12$.
\end{proposition}


\begin{proof}
\begin{enumerate}[label=\arabic*.,leftmargin=*]
\item We apply \citet[Theorem 1.1 and Corollary 1.2]{demarco2017}.
Their hypotheses are a number field, a smooth projective
base curve, an elliptic surface and a rational section of
nonzero generic canonical height. Here they are $\Q$,
$\mathbb P^1$, the displayed rational elliptic surface,
and $P$, whose generic height for $(O)$ is $1/12$ by \ref{ari:C.6.4}.
At a torsion specialization the canonical height is zero.
Corollary 1.2 consequently gives convergence of every
nonrepeating sequence of Galois-orbit probability measures.

\item The limiting measure is the curvature of the height metric
divided by the degree of its line bundle. In our normalization,
\ref{ari:C.8.9} computes the curvature as $m_P$, and \ref{ari:C.6.4} computes
that degree as $1/12$. Hence the probability measure is
$12m_P$. This division by the degree is essential: a measure
of mass $1/12$ cannot be a weak limit of probability measures.
The source's informal description of the pullback current
is interpreted with this normalization.

\item For completeness, the passage from individual orbits to
the full division divisor does not require uniformity of
the equidistribution theorem in $n$. The surface height
formula, with our half-self-pairing convention, reads
\[
 \widehat h_O(Q)=\chi(\mathcal O_{\mathcal E})+(Q\cdot O)
                -\frac12\sum_v\operatorname{contr}_v(Q).
\]
Here $\chi=1$, and the finitely many fibre corrections are
bounded as $Q=nP$ varies, because they depend only on its
component. Since $\widehat h_O(nP)=n^2/12$, it follows
that $(nP\cdot O)=n^2/12+O(1)$. Removing the singular fibres
changes this by $O(1)$: the local intersection at each
such fibre is bounded as $n$ varies in characteristic zero.
One can see this in the formal group after a fixed
multiple enters it; multiplication by a nonzero integer
is a local analytic isomorphism over the residue field
$\C$. Thus $\deg D_n=n^2/12+O(1)$.

\item At any fixed smooth torsion parameter $t_0$, of exact order
$d$, the multiplicity in $D_n$ is zero unless $d\mid n$,
and is constant when $d\mid n$. Indeed the formal parameter
of $nP=[n/d](dP)$ has first term $(n/d)$ times that of $dP$,
with nonzero constant coefficient. Fix a continuous test
function and $\varepsilon>0$. Corollary 1.2 implies that
all but finitely many distinct Galois orbits of torsion
parameters have average within $\varepsilon$ of the limiting
integral; otherwise a countersequence would violate it.
The total weighted contribution of the finitely many
exceptional orbits is bounded independently of $n$, by
the multiplicity observation. Divide by $\deg D_n$ and
let $n\to\infty$. This proves the probability convergence
for $D_n$, and the degree asymptotic proves the unnormalized
version.
\end{enumerate}
\end{proof}

\begin{proposition}[the two first-derivative identities]\label{ari:C.8.11}
\StatementRecord{ari:C.8.11}{proved}{Historical events: Claude Fable 5.2 (Anthropic), effort high, 16 November 2026: used and rechecked in this run. GPT-6 Astra (OpenAI), ChatGPT Pro mode, effort Pro, S6\_09\_PAPERS, 17 November 2026: targeted proof check. Also checked in this writing. Pre-continuation writing events: GPT-6 Astra (OpenAI), ChatGPT Pro mode, effort Pro, S6\_10\_PAPERS, 18 November 2026: targeted proof check. Continuation metadata events: GPT-6 Astra (OpenAI), ChatGPT Pro mode, effort Pro, S6\_10\_PAPERS continuation, 19 November 2026: used and rechecked in this run; status-vocabulary reconciliation; metadata reconciliation only; no new mathematical check. \emph{Pre-continuation recorded proof scope.} Read complete proof and re-derived its Green-function divisor/mean, beta integral, signed Gaussian Fourier, Betti Jacobian/current-density, division-divisor multiplicity, or heat-equation calculation. For C.8.10 the DeMarco-Mavraki external theorem locator remains a separate gate.}{the hypotheses, notation and definitions stated here; the explicit constructions and calculations in the proof}

On an ordinary modulus chart put
\[
 a=\frac{\Im\mu}{T},\qquad
 W=\frac{\vartheta_1'(\mu,\tau)}{\vartheta_1(\mu,\tau)}
                         +2\pi ia,\qquad
 R=\frac{d\mu}{d\tau}-a.
\]
Then
\[
 L_\tau=\frac{W}{8\pi i}(W+4\pi iR),\qquad
 \det\frac{d}{d\tau}
 \begin{pmatrix}\tau&\mu\\\mu&\beta^{\rm part}/6\end{pmatrix}
       =\frac{(W+2\pi iR)^2}{4\pi^2}.
\]
Moreover
\[
 R=\frac{g_2}{6\pi i\wp'(\mu)}-\frac{W}{2\pi i},
 \qquad
 L_\tau=\frac{iW(3W\wp'(\mu)-2g_2)}
                         {24\pi\wp'(\mu)}.
\]
Multiplication by $6/\pi$ gives the critical-point factorization for
$D_{\rm part}$ in Proposition~\ref{push:src:D.4.4}.
\end{proposition}


\begin{proof}
\begin{enumerate}[label=\arabic*.,leftmargin=*]
\item Put $f=\vartheta_1'/\vartheta_1$ and
$E=\pi^2E_2/3$. The theta heat equation and
$\eta'/\eta=\pi iE_2/12$ give
\[
 \partial_\tau\log(\vartheta_1/\eta)
       =\frac{f_z+f^2+E}{4\pi i}.
\]
At $\wp(\mu)=0$, the identity $f_z=-\wp-E$ reduces
this to $f^2/(4\pi i)$. Thus, along the section,
\[
 \frac{d}{d\tau}\log(\vartheta_1(\mu,\tau)/\eta)
       =f\mu'+\frac{f^2}{4\pi i}.
\]
Differentiate the Gaussian $-\pi(\Im\mu)^2/T$ using
Wirtinger derivatives, retaining the factor $1/2$ in
the derivative of the real part of a holomorphic logarithm.
Substitution of $f=W-2\pi ia$ and $\mu'=R+a$ gives the
first formula. Since
$\beta^{\rm part}/6=(i/\pi)\log(\vartheta_1/(2\sqrt3\eta))$,
its derivative is $if\mu'/\pi+f^2/(4\pi^2)$.
Subtracting $\mu'^2$ gives
$(f+2\pi i\mu')^2/(4\pi^2)$, the second formula.

\item Finally differentiate $\wp_\tau(\mu(\tau))=0$ by the theta
heat equation. After substituting $\wp(\mu)=0$,
$\wp''(\mu)=-g_2/2$, and $f=W-2\pi ia$, the result is
$6\pi i\wp'R=g_2-3W\wp'$. This is the formula for $R$.
Inserting it into the first identity gives the last one.
These computations independently verify the differential formula used in Proposition~\ref{ari:C.3.2};
its apparent singularity at a half-period is
handled by the root chart rather than division by zero.
\end{enumerate}
\end{proof}



\section{Critical points and the global level curves}

We continue to write $L(t)=\log N(\mu(\tau),\tau)$, where
$t=j(\tau)/1728$ and $\wp_\tau(\mu)=0$. The sign and lattice
translation of $\mu$ do not affect any expression below. We use
Eichler and Zagier's unnumbered theorem on p.~399 of
\cite{EichlerZagier1982}, also stated as equation (6) and the
unnumbered theorem in Section~2 of \cite{DukeImamoglu2008}.
For the cusp branch in this problem their theorem says
\begin{equation}\label{ari:eq:cc-EZ}
 \begin{aligned}
 \mu&=\frac12-i\beta_{\rm cusp}-
       \frac{36i\sqrt6}{\pi}\sum_{n\ge1}\frac{a_nq^n}{n^2},
 &\beta_{\rm cusp}&=\frac{\log(5+2\sqrt6)}{2\pi},\\
 \sum_{n\ge1}a_nq^n&=\frac{\Delta_0}{E_6^{3/2}},
 &\mu''&=144\pi i\sqrt6\frac{\Delta_0}{E_6^{3/2}}.
 \end{aligned}
\end{equation}
The branch of $E_6^{3/2}$ tends to one at the cusp. Thus the
formula applies directly on $\Im\tau>1$ and then by analytic
continuation away from the half-period branch points. Our use at $i$ is on the square-root chart specified explicitly below.
The original theorem and the hypergeometric reformulation concern
the logarithm $\mu$; neither asserts the global inequality for
$L$ proved here. Searches through these papers and their cited
work, and through the torus Green-function literature, did not
locate the radial monotonicity statement below.

\begin{proposition}[a radial derivative identity]\label{ari:C.8.12}
\StatementRecord{ari:C.8.12}{proved}{Historical events: Claude Fable 5.2 (Anthropic), effort high, 15 November 2026: carried as background without fault found. GPT-6 Astra (OpenAI), ChatGPT Pro mode, effort Pro, S6\_09\_PAPERS, 17 November 2026: targeted proof check. Also checked in this writing. Pre-continuation writing events: GPT-6 Astra (OpenAI), ChatGPT Pro mode, effort Pro, S6\_10\_PAPERS, 18 November 2026: targeted proof check. Continuation metadata events: GPT-6 Astra (OpenAI), ChatGPT Pro mode, effort Pro, S6\_10\_PAPERS continuation, 19 November 2026: used and rechecked in this run; status-vocabulary reconciliation; metadata reconciliation only; no new mathematical check. \emph{Pre-continuation recorded proof scope.} Read complete proof. Re-derived critical factor, radial derivative, cusp Fourier bound, Cayley ramification removal, rational-series tail bounds, scalar Bernstein implication and raywise Jordan-curve argument. The precise imported elliptic-function formula remains an external-locator gate until opened.}{the hypotheses, notation and definitions stated here; \cref{ari:C.8.11}; \cref{ari:C.3.2}}

Put $p=\wp'_\tau(\mu)$ and
\[
 X=\frac{3Wp}{g_2},\qquad F=1-X,
 \qquad W=\frac{\vartheta_1'(\mu,\tau)}{\vartheta_1(\mu,\tau)}
              +\frac{2\pi i\Im\mu}{\Im\tau}.
\]
These are single-valued modular invariants wherever $g_2\ne0$.
They have continuous extensions at $t=0,1$, with
$X(0)=1/4$ and $X(1)=0$. At every $t\ne0,1$,
\begin{equation}\label{ari:eq:cc-radial}
 (t-1)L_t=\frac{X(X-2)}{24},\qquad
 r\frac{d}{dr}L(1+re^{i\theta})
       =\frac{\Re(X(X-2))}{12}.
\end{equation}
Also, if $a=\Im\mu/\Im\tau$ and $R=\mu'-a$, then
\begin{equation}\label{ari:eq:cc-FR}
 F=\frac{6\pi ip}{g_2}R.
\end{equation}
Consequently $|F|<1$ implies strict decrease along the ray
from $1$ through $t$.
\end{proposition}


\begin{proof}
\begin{enumerate}[label=\arabic*.,leftmargin=*]
\item A lattice translation leaves $W$ unchanged, and replacing $\mu$
by $-\mu$ changes the signs of both $W$ and $p$. Under a modular
change, $W,p,g_2$ have respective weights $1,3,4$, so $X$ is
invariant. The two identities of \ref{ari:C.8.11} give
\[
 L_\tau=\frac{iW(3Wp-2g_2)}{24\pi p},\qquad
 R=\frac{g_2-3Wp}{6\pi ip}.
\]
Ramanujan's identities and $p^2=-g_3$ give
\[
 t'_\tau=-2\pi it\frac{E_6}{E_4},\qquad
 \frac{t-1}{t'_\tau}=-\frac{E_6}{2\pi iE_4}.
\]
Substitution, using $g_2=4\pi^4E_4/3$ and
$g_3=8\pi^6E_6/27$, proves \eqref{ari:eq:cc-radial} and
\eqref{ari:eq:cc-FR}. Since $X(X-2)=F^2-1$, $|F|<1$ makes
its real part strictly negative.

\item At $i$, $W=p=0$ and $g_2\ne0$, so $X\to0$. At $\rho$, put $\delta=\tau-\rho$ and $p_0=\wp'_\rho(\mu_0)$.
The expansion in \ref{ari:C.3.2} and Ramanujan's identities give
\[
 R=\frac{3p_0}{8\pi^2}\delta+O(|\delta|^2),\qquad
 g_2=\frac{3ip_0^2}{\pi}\delta+O(\delta^2).
\]
Substitution into \eqref{ari:eq:cc-FR} gives $F(\rho)=3/4$.
Equivalently, the two possible limits furnished by
\eqref{ari:eq:cc-radial} and $L_t(0)=7/384$ are $X=1/4,7/4$;
the expansion of $R$ selects $1/4$. The same expansions show
continuity on the coarse $t$-line.
\end{enumerate}
\end{proof}

\begin{proposition}[the cusp exclusion]\label{ari:C.8.13}
\StatementRecord{ari:C.8.13}{proved}{Historical events: GPT-6 Astra (OpenAI), ChatGPT Pro mode, effort Pro, S6\_09\_PAPERS, 17 November 2026: targeted proof check. Also checked in this writing. Pre-continuation writing events: GPT-6 Astra (OpenAI), ChatGPT Pro mode, effort Pro, S6\_10\_PAPERS, 18 November 2026: computation or formal certificate; targeted proof check. Continuation metadata events: GPT-6 Astra (OpenAI), ChatGPT Pro mode, effort Pro, S6\_10\_PAPERS continuation, 19 November 2026: used and rechecked in this run; status-vocabulary reconciliation; metadata reconciliation only; no new mathematical check. \emph{Pre-continuation recorded proof scope.} Read complete proof. Re-derived critical factor, radial derivative, cusp Fourier bound, Cayley ramification removal, rational-series tail bounds, scalar Bernstein implication and raywise Jordan-curve argument. The precise imported elliptic-function formula remains an external-locator gate until opened. Reran Ccritical\_rational\_certificate.py and Ccritical\_scalar.py. All 100 exact ODE coefficients and outward integer bounds were recomputed; all 15 exact Q(sqrt(3)) Bernstein coefficients are positive. Saved stdout and exact output certificates accompany the code.}{the hypotheses, notation and definitions stated here; \cref{ari:eq:cc-EZ}; \cref{ari:eq:ca-exp}}

If $\Im\tau\ge7/5$, then $|F|<17/20$.
\end{proposition}


\begin{proof}
\begin{enumerate}[label=\arabic*.,leftmargin=*]
\item Here $|q|<1/6600$. On the auxiliary circle $|q|=1/2000$,
the Lambert series and eta product give
\[
 |E_6-1|\le
 504r\frac{1+26r+66r^2+26r^3+r^4}{(1-r)^7},\qquad
 |\Delta_0|\le r\exp\!\frac{24r}{(1-r)^2},
 \quad r=\frac1{2000}.
\]
The right sides imply $|\Delta_0/E_6^{3/2}|<1/1250$.
In particular $E_6$ has no zero on this disc and the specified
branch is holomorphic there. By Cauchy's inequality, with
$s=10/33$,
\[
 \sum_{n\ge1}\frac{|a_nq^n|}{n}
 \le\frac{-\log(1-s)}{1250},\qquad
 \sum_{n\ge1}\frac{|a_nq^n|}{n^2}
 \le\frac1{1250}\left(s+\frac{s^2}{4(1-s)}\right).
\]
Equation \eqref{ari:eq:cc-EZ} therefore yields
\[
 |\mu-(1/2-i\beta_{\rm cusp})|<\frac1{125},\qquad
 |\mu'|<\frac{13}{250},\qquad \beta_{\rm cusp}<\frac{365}{1000}.
\]
For $r=1/6600$, the corresponding weight-four and weight-six
Lambert bounds are
\[
 240r\frac{1+4r+r^2}{(1-r)^5}<\frac{37}{1000},\qquad
 504r\frac{1+26r+66r^2+26r^3+r^4}{(1-r)^7}
       <\frac{77}{1000}.
\]
Thus $|E_4|>963/1000$ and $|E_6|^{1/2}<1038/1000$.
Since $6\pi|p/g_2|=\sqrt6|E_6|^{1/2}/|E_4|$,
\begin{equation}\label{ari:eq:cc-cuspbound}
 |F|<\sqrt6\frac{1038}{963}
 \left\{\frac{365/1000+1/125}{7/5}+\frac{13}{250}\right\}
 <\frac{17}{20}.
\end{equation}
All scalar comparisons here follow from geometric series and
the exponential Taylor bounds \eqref{ari:eq:ca-exp}; for the logarithm
use $-\log(1-s)=\sum s^n/n$ with its geometric remainder.
For reference, the upper quantities before the last rounding are
less than $789/10^6$, $755/10^5$, $51/1000$, and $841/1000$
for $|\Delta_0/E_6^{3/2}|$, the error in $\mu$, the bound on
$\mu'$, and the final bound on $|F|$, respectively.
\end{enumerate}
\end{proof}

\begin{proposition}[the compact exclusion by one Cayley disc]\label{ari:C.8.14}
\StatementRecord{ari:C.8.14}{proved}{Historical events: GPT-6 Astra (OpenAI), ChatGPT Pro mode, effort Pro, S6\_09\_PAPERS, 17 November 2026: targeted proof check. Also checked in this writing. Pre-continuation writing events: GPT-6 Astra (OpenAI), ChatGPT Pro mode, effort Pro, S6\_10\_PAPERS, 18 November 2026: computation or formal certificate; targeted proof check. Continuation metadata events: GPT-6 Astra (OpenAI), ChatGPT Pro mode, effort Pro, S6\_10\_PAPERS continuation, 19 November 2026: used and rechecked in this run; status-vocabulary reconciliation; metadata reconciliation only; no new mathematical check. \emph{Pre-continuation recorded proof scope.} Read complete proof. Re-derived critical factor, radial derivative, cusp Fourier bound, Cayley ramification removal, rational-series tail bounds, scalar Bernstein implication and raywise Jordan-curve argument. The precise imported elliptic-function formula remains an external-locator gate until opened. PASS: unchanged Ccritical\_rational\_certificate.py recomputed exact modular ODE coefficients and outward-rounded rational coefficient norms; Ccritical\_scalar.py recomputed all 15 positive Bernstein coefficients over Q(sqrt(3)). Outputs code/I/replay\_arithmetic/. Reran Ccritical\_rational\_certificate.py and Ccritical\_scalar.py. All 100 exact ODE coefficients and outward integer bounds were recomputed; all 15 exact Q(sqrt(3)) Bernstein coefficients are positive. Saved stdout and exact output certificates accompany the code.}{the hypotheses, notation and definitions stated here; \cref{ari:eq:cc-EZ}; \cref{ari:eq:cc-Fseries}; \cref{ari:eq:cc-ode}; \cref{ari:eq:cc-recurrence}; \cref{ari:eq:cc-power}; \cref{ari:eq:cc-finite-table}; \cref{ari:eq:cc-norms}; \cref{ari:eq:cc-Ffactored}; \cref{ari:eq:cc-scalar}; \cref{ari:C.8.12}}

Let $\chi=(\tau-i)/(\tau+i)$. On
$0<|\chi|\le2-\sqrt3$, the invariant $F$ satisfies $|F|<1$.
At the two equianharmonic points on the boundary its continuous
value is $3/4$.
\end{proposition}


\begin{proof}
\begin{enumerate}[label=\arabic*.,leftmargin=*]
\item Set
\[
 \begin{gathered}
 r_0=2-\sqrt3,\qquad A=E_4(i),\qquad k=\pi^2A,\\
 v=\pi\sqrt A\,\chi,\qquad z=v^2=k\chi^2,\\
 z_0=kr_0^2,\qquad c_0=1-\frac1{\sqrt3}.
 \end{gathered}
\]
In this proof $z$ is a coordinate on the modular disc, not a
point of the elliptic curve. Define analytic functions $f_2,f_4,f_6$
by
\[
 \begin{gathered}
 \sqrt A f_2(v)=\frac{E_2(\tau)}{(1-\chi)^2}
                         -\frac3{\pi(1-\chi)},\\
 A f_4(v)=\frac{E_4(\tau)}{(1-\chi)^4},\qquad
 A^{3/2} f_6(v)=\frac{E_6(\tau)}{(1-\chi)^6}.
 \end{gathered}
\]
Ramanujan's differential identities become the rational system
\begin{equation}\label{ari:eq:cc-ode}
 \begin{gathered}
 f_2'=\frac{f_4-f_2^2}{3},\qquad
 f_4'=\frac43(f_6-f_2f_4),\qquad
 f_6'=2(f_4^2-f_2f_6),\\
 (f_2,f_4,f_6)(0)=(0,1,0).
 \end{gathered}
\end{equation}
Uniqueness makes $f_4$ even and $f_2,f_6$ odd. Write
\[
 B(z)=f_4(v),\quad J(z)=\frac{f_6(v)}{2v},\quad
 D(z)=\frac{B(z)^3-4zJ(z)^2}{J(z)^{3/2}},\qquad
 H(z)=\sum_{n\ge0}\frac{-[z^n]D(z)}{16n^2-1}z^n.
\]
The branch has $J(0)^{1/2}=1$. The estimates below establish
its regularity on a larger disc, so these definitions are
unambiguous on $|z|\le z_0$.

\item Write $h=(1-\chi)(\mu-(1-\tau)/2)$. The Möbius chain rule and
\eqref{ari:eq:cc-EZ} give
\[
 h''(\chi)=-\frac{4\mu''(\tau)}{(1-\chi)^3}
       =-\frac{i}{2\sqrt{3\pi}}\chi^{-3/2}D(k\chi^2).
\]
The root expansion at $i$ gives the branch
\begin{equation}\label{ari:eq:cc-h}
 h=\frac{2i}{\sqrt{3\pi}}\sqrt \chi\,H(z).
\end{equation}
Indeed differentiating this series twice gives precisely the
coefficient denominator $16n^2-1$ above. The two possible affine
homogeneous terms are absent because the half-period difference
has a Puiseux expansion consisting only of half-integral powers.
Changing the sign of the square root changes the signs of $h,p,R$
and leaves $F$ unchanged.

\item Put $r=|\chi|$. The definitions of $a$ and $R$ give exactly
\[
 R=\frac{1-\chi}{2i}
       \left(h'+\frac{\overline \chi h+\overline h}{1-r^2}\right),
 \qquad
 \frac p{g_2}=-\frac{i\sqrt \chi}{\sqrt{3\pi}(1-\chi)}
                         \frac{\sqrt J}{B}.
\]
Consequently
\begin{equation}\label{ari:eq:cc-Fseries}
 F=\frac{\sqrt J}{B}
 \left\{H+4zH'+\frac{2r}{1-r^2}(rH-\overline H)\right\}.
\end{equation}
All conjugates in this display are values at the same point;
$H$ has real coefficients, so $\overline{H(z)}=H(\overline z)$.

\item The apparent pole at the equianharmonic points must be removed
before estimating this expression. We have $B(-z_0)=0$.
At these points $R=0$, since the expansion at $\rho$ gives
$\mu'=-1/3=a$. Thus, if
\[
 S=\sqrt J H,\qquad T=\sqrt J(H+4zH'),
\]
then $T(-z_0)=c_0S(-z_0)$. Define the analytic quotients
\[
 C(z)=\frac{B(z)}{1+z/z_0},\qquad
 U(z)=\frac{T(z)-c_0S(z)}{(1+z/z_0)C(z)},\qquad
 V(z)=\frac{\sqrt{J(z)}}{C(z)}.
\]
The estimates below prove $C\ne0$ on the needed larger disc.
For $r<r_0$, equation \eqref{ari:eq:cc-Fseries} is now
\begin{equation}\label{ari:eq:cc-Ffactored}
 F=U+V\frac{dH+q(H-\overline H)}{1+z/z_0},
 \qquad
 d=\frac{2(r_0-r)}{(1+r_0)(1+r)},\quad q=\frac{2r}{1-r^2}.
\end{equation}
Here $q$ is only a real abbreviation within this display and
its estimates, not the Fourier nome.

\item We prove the following four coefficient bounds. For an analytic
power series $K=\sum K_nz^n$, put
$\|K\|_s=\sum|K_n|s^n$. Then, for $0\le s\le z_0$,
\begin{equation}\label{ari:eq:cc-norms}
 \begin{aligned}
 \|U\|_s&\le\frac1{\sqrt3}+\frac{13}{200}s,&
 \|V\|_s&\le1+\frac7{50}s,\\
 \|H\|_s&\le1+\frac{113}{1000}s,&
 \|H'\|_s&\le\frac{13}{100}.
 \end{aligned}
\end{equation}
We give the coefficient arithmetic and all tails, rather than
leaving these inequalities to sampling a circle. In
\eqref{ari:eq:cc-ode}, write $f_2=\sum a_nv^n$,
$f_4=\sum b_nv^n$, $f_6=\sum c_nv^n$. Starting with
$(a_0,b_0,c_0)=(0,1,0)$ (only in this recurrence the symbol
$c_0$ denotes the zero initial coefficient), compute
\begin{equation}\label{ari:eq:cc-recurrence}
 \begin{aligned}
 a_{n+1}&=\frac{b_n-\sum_{j=0}^na_ja_{n-j}}{3(n+1)},\\
 b_{n+1}&=\frac{4(c_n-\sum_{j=0}^na_jb_{n-j})}{3(n+1)},\\
 c_{n+1}&=\frac{2(\sum_{j=0}^nb_jb_{n-j}-
                          \sum_{j=0}^na_jc_{n-j})}{n+1}.
 \end{aligned}
\end{equation}
Thus $B_n=b_{2n}$ and $J_n=c_{2n+1}/2$. For powers of a series
$K$ with $K_0=1$, the equally elementary recurrence is
\begin{equation}\label{ari:eq:cc-power}
 [z^n]K^\alpha=\frac1n\sum_{j=1}^n
                 ((\alpha+1)j-n)K_j[z^{n-j}]K^\alpha.
\end{equation}
Together with $H_n=-D_n/(16n^2-1)$, these determine every
coefficient used. For example,
\[
 \begin{array}{c|rrrr}
 n&0&1&2&3\\ \hline
 B_n&1&10/9&5/27&4/81\\
 J_n&1&14/27&28/135&14/405\\
 H_n&1&13/135&-1999/153090&29299/28146690.
 \end{array}
\]
Division by $1+z/z_0$ uses $C_0=1$ and
$C_n=B_n-C_{n-1}/z_0$; the same recurrence applies to
$(T-c_0S)/(1+z/z_0)$. Division by $C$ uses the usual
convolution recurrence with constant coefficient one.

\item Here are sufficient rational enclosures for the only nonrational
constants in this arithmetic:
\begin{equation}\label{ari:eq:cc-constants}
 \begin{gathered}
 14<k<72/5,\qquad
 1.0315619043558156493568263343177356527388253980203
 <z_0\\
 <1.0315619043558156493568263343177356527388253980205.
 \end{gathered}
\end{equation}
They follow without special-function evaluation from
$\pi=16\arctan(1/5)-4\arctan(1/239)$, the first 100 terms
of the alternating arctangent series, the exponential series
through degree 160, and
\[
 A=1+240\sum_{n=1}^{49}\frac{n^3e^{-2\pi n}}{1-e^{-2\pi n}}
       +\varepsilon_A,\qquad
 0<\varepsilon_A<
 \frac{240(4e^{-2\pi})^{50}}
 {(1-4e^{-2\pi})(1-e^{-2\pi})}.
\]
The last bound uses $n^3\le4^n$. Bounds for $\sqrt3$ follow
by squaring rational endpoints. All ensuing coefficient
comparisons are rational arithmetic using these enclosures.

\item To bound infinite tails, first consider $|z|=12/5$.
Here $|\chi|<5/12$, $\Im\tau>7/17$ and
$|e^{2\pi i\tau}|<1/13$. Lambert series give
\begin{equation}\label{ari:eq:cc-outer}
 |B|<400,\qquad |J|<3000\qquad (|z|=12/5).
\end{equation}
For clarity the respective upper bounds before rounding are
\[
 \left(\frac{12}{7}\right)^4
 \left(1+240u\frac{1+4u+u^2}{(1-u)^5}\right)<400,
 \quad u=1/13,
\]
\[
 \frac13\left(\frac{12}{7}\right)^6
 \left(1+504u\frac{1+26u+66u^2+26u^3+u^4}{(1-u)^7}\right)<3000.
\]
We have used $A>1$ and $2\sqrt{12/5}>3$.
For $R=3/2$, let $\xi=5/8$ and retain coefficients through
degree 59. The tails of $B,J,C$ are respectively at most
\begin{equation}\label{ari:eq:cc-tail60}
 \frac{400\xi^{60}}{1-\xi},\qquad
 \frac{3000\xi^{60}}{1-\xi},\qquad
 \frac{400z_0}{12/5-z_0}\frac{\xi^{60}}{1-\xi}.
\end{equation}
For the last bound, $B(-z_0)=0$ implies
$|C_n|\le400z_0(12/5)^{-n}/(12/5-z_0)$.
The finite sums from \eqref{ari:eq:cc-recurrence}--\eqref{ari:eq:cc-power}
with these tails give
\[
 \begin{array}{c|cccc}
 &\|B\|_{3/2}&\|J\|_{3/2}&\|C-1\|_{3/2}
                   &\min_{|z|\le3/2}\Re J(z)\\ \hline
 \text{bound}&<3276/1000&<2411/1000&<335/1000&>2049/10000.
 \end{array}
\]
For the last column, the lower bound used is exactly
\[
 1-J_2R^2-\frac{J_1^2}{8J_2}
       -\sum_{n=3}^{59}|J_n|R^n-
                   \frac{3000\xi^{60}}{1-\xi}.
\]
This is obtained by minimizing the quadratic
$\Re(1+J_1z+J_2z^2)$ on $|z|=R$; the minimum occurs at
$\cos\arg z=-J_1/(4J_2R)$. In particular, on $|z|\le3/2$,
\[
 |B|<33/10,\quad |J|<5/2,\quad |C|>3/5,\quad
 \Re J>1/5,\quad
 |D|\le\frac{(33/10)^3+6(5/2)^2}{(1/5)^{3/2}}<1000.
\]
Cauchy's inequality and the denominators in $H_n$ now give
\[
 |H|\le1+1000\sum_{n\ge1}\frac1{16n^2-1}
       =1+1000\left(\frac12-\frac\pi8\right)<110
                      \quad (|z|\le3/2).
\]
On $|z|\le7/5$, the same argument gives
\[
 |H+4zH'|\le1+1000\sum_{n\ge1}
                 \frac{(14/15)^n}{4n-1}
       \le1+\frac{1000}{3}\log15<910.
\]
It follows on $|z|=7/5$ that $|U|<10000$ and $|V|<3$;
indeed $|(1+z/z_0)C|>(7/(5z_0)-1)3/5$, and
$|\sqrt J|<8/5$.

\item Finally retain coefficients through degree 99 and put
$\eta=5z_0/7$. The tails in the coefficient norms for
$U,V,H,H'$ at $z_0$ are bounded respectively by
\begin{equation}\label{ari:eq:cc-tail100}
 \frac{10000\eta^{100}}{1-\eta},\quad
 \frac{3\eta^{100}}{1-\eta},\quad
 \frac{110\eta^{100}}{1-\eta},\quad
 \frac{110\eta^{100}}{z_0}
 \left(\frac{100}{1-\eta}+\frac{\eta}{(1-\eta)^2}\right).
\end{equation}
The finite coefficient sums plus precisely these tails give
\begin{equation}\label{ari:eq:cc-finite-table}
 \begin{array}{c|cccc}
 & (\|U\|_{z_0}-1/\sqrt3)/z_0
 & (\|V\|_{z_0}-1)/z_0
 & (\|H\|_{z_0}-1)/z_0 & \|H'\|_{z_0}\\ \hline
 \text{upper bound}&635/10000&1368/10000&1115/10000&1293/10000.
 \end{array}
\end{equation}
The recurrences, cutoffs and remainders above specify each finite
comparison completely. The attached exact integer/rational
script reproduces them with outward rounding on denominator
$10^{100}$ and independently checks every rational ODE
coefficient; no point of the modular domain is enumerated.
Equations \eqref{ari:eq:cc-finite-table} imply
\eqref{ari:eq:cc-norms}, with room in each inequality.

\item We finish with a single polynomial inequality. For $s=|z|=kr^2$,
\[
 |H-\overline H|\le2|\Im z|\,\|H'\|_s,\qquad
 \frac{|\Im z|}{|1+z/z_0|}\le s.
\]
The second inequality follows by writing $z/z_0=ue^{i\phi}$
and using $(\cos\phi+u)^2\ge0$. The useful cancellation is
\[
 |dH+q(H-\overline H)|^2
      =d^2|H|^2+4q(d+q)(\Im H)^2.
\]
Also
\[
 \frac d{|1+z/z_0|}\le
 e(r):=\frac{c_0}{(1+r)(1+r/r_0)},\qquad
 q\le\kappa r,\quad d+q\le c_0+\kappa r^2,
 \quad\kappa=1+\frac2{\sqrt3}.
\]
The last bounds use $(1-r_0^2)^{-1}=1/2+1/\sqrt3$.
It follows from \eqref{ari:eq:cc-Ffactored}--\eqref{ari:eq:cc-norms},
using $k<72/5$, that
\begin{equation}\label{ari:eq:cc-scalar}
 |F|\le U_0(r)+V_0(r)
 \sqrt{e(r)^2H_0(r)^2+
       4\kappa r(c_0+\kappa r^2)(72/5)^2r^4(13/100)^2},
\end{equation}
where
\[
 U_0=\frac1{\sqrt3}+\frac{13}{200}\frac{72}{5}r^2,
 \quad V_0=1+\frac7{50}\frac{72}{5}r^2,
 \quad H_0=1+\frac{113}{1000}\frac{72}{5}r^2.
\]
To verify that the right side is less than one, let
\[
 M(r)=(1+r)^2(1+r/r_0)^2,
\]
\begin{equation}\label{ari:eq:cc-P}
 \begin{aligned}
 P(r)={}&(1-U_0)^2M(r)\\
 &-V_0^2\left\{c_0^2H_0^2+
 4\kappa r(c_0+\kappa r^2)(72/5)^2r^4(13/100)^2M(r)\right\}.
 \end{aligned}
\end{equation}
This polynomial has $P(0)=0$. Write
\[
 \frac{P(r_0x)}{r_0x}
     =\sum_{j=0}^{14}b_j\binom{14}{j}x^j(1-x)^{14-j}.
\]
All coefficients belong to $\Q(\sqrt3)$. Explicit lower bounds,
obtained by expanding the displayed factored polynomial, are
\[
 \begin{array}{c|rrrrrrrr}
 j&0&1&2&3&4&5&6&7\\ \hline
 100b_j>&169&175&181&187&192&197&202&205\\[2pt]
 j&8&9&10&11&12&13&14&\\ \hline
 100b_j>&207&207&204&197&184&160&118&.
 \end{array}
\]
For an exact independent recipe, if $P(r)/r=\sum q_nr^n$,
then
\[
 b_j=\sum_{n=0}^j q_nr_0^n
                  \frac{\binom jn}{\binom{14}n}.
\]
For example $b_0=4-4\sqrt3/3$ and
$b_1=25357/13125-453\sqrt3/4375$. The rational enclosure
\[
 \begin{gathered}
 1.73205080756887729352744634150587236694280525381038<\sqrt3,\\
 \sqrt3<1.73205080756887729352744634150587236694280525381039.
 \end{gathered}
\]
verifies every entry by integer multiplication; its endpoints
are checked by squaring. The complete exact coefficient list is
also retained with the code. Every Bernstein basis function is
nonnegative on $[0,1]$, so $P(r)>0$ for $0<r\le r_0$.
Since $U_0<1$ there, squaring \eqref{ari:eq:cc-scalar} is legitimate
and proves $|F|<1$. Formula \eqref{ari:eq:cc-Ffactored} applies
at $r=r_0$ except at $z=-z_0$ by continuity; at that remaining
point \ref{ari:C.8.12} gives $F=3/4$. This proves the assertion.
\end{enumerate}
\end{proof}

\begin{proposition}[every critical point and every level set]\label{ari:C.8.15}
\StatementRecord{ari:C.8.15}{proved}{Historical events: Claude Fable 5.2 (Anthropic), effort high, 15 November 2026: carried as background without fault found. Claude Fable 5.2 (Anthropic), effort high, 16 November 2026: used and rechecked in this run. GPT-6 Astra (OpenAI), ChatGPT Pro mode, effort Pro, S6\_09\_PAPERS, 17 November 2026: targeted proof check. Also checked in this writing. Pre-continuation writing events: GPT-6 Astra (OpenAI), ChatGPT Pro mode, effort Pro, S6\_10\_PAPERS, 18 November 2026: computation or formal certificate; targeted proof check. Continuation metadata events: GPT-6 Astra (OpenAI), ChatGPT Pro mode, effort Pro, S6\_10\_PAPERS continuation, 19 November 2026: used and rechecked in this run; status-vocabulary reconciliation; metadata reconciliation only; no new mathematical check. \emph{Pre-continuation recorded proof scope.} Read complete proof. Re-derived critical factor, radial derivative, cusp Fourier bound, Cayley ramification removal, rational-series tail bounds, scalar Bernstein implication and raywise Jordan-curve argument. The precise imported elliptic-function formula remains an external-locator gate until opened. Reran Ccritical\_rational\_certificate.py and Ccritical\_scalar.py. All 100 exact ODE coefficients and outward integer bounds were recomputed; all 15 exact Q(sqrt(3)) Bernstein coefficients are positive. Saved stdout and exact output certificates accompany the code.}{the hypotheses, notation and definitions stated here; \cref{ari:C.8.13}; \cref{ari:C.8.14}; \cref{ari:eq:cc-radial}; \cref{ari:eq:cc-EZ}; \cref{fam:src:C.2.3}}

For every zero $\mu$ of $\wp_\tau$, the zero sets of the two
critical factors are exactly
\[
 \begin{aligned}
 W=0&\quad\Longleftrightarrow\quad
       \tau\in\operatorname{SL}_2(\Z)\{i,\rho\},\\
 3W\wp'_\tau(\mu)-2g_2=0&\quad\Longleftrightarrow\quad
       \tau\in\operatorname{SL}_2(\Z)\rho.
 \end{aligned}
\]
On the coarse $t$-plane the function $L$ has no critical point
away from $t=1$. The point $t=0$ is a noncritical $C^1$ point,
with $L_t(0)=7/384$. The point $t=1$ is the unique global
maximum, of value $\log\sqrt2$, and has the fractional-power
local maximum described in C.3. In particular there is no
saddle in $0<t<1$. On the root charts, $h=0$ over $t=1$
is a nondegenerate maximum, while $\delta=0$ over $t=0$
is a degenerate critical point with nonzero harmonic cubic
leading term and three alternating ascending and descending
sectors. These are all critical points on the root charts.

For every angle $\theta$, $r\mapsto L(1+re^{i\theta})$ is
strictly decreasing for $r>0$, from $\log\sqrt2$ to $-\infty$.
For each $c<\log\sqrt2$, the level set $\{L=c\}$ is one
Jordan curve, strictly star-shaped about $1$. Its bounded
complement is $\{L>c\}$; its unbounded complement is
$\{L<c\}$. The curve is real analytic except possibly at
$t=0$ when $c=\frac14\log3$, where it is $C^1$ with nonzero
tangent. The closed superlevel sets form nested topological
discs. At the maximum the level set is the single point $1$,
and above it the level set is empty. The corresponding
statements for $D_{\rm part}$ follow by its increasing affine
change from $L$.
\end{proposition}


\begin{proof}
\begin{enumerate}[label=\arabic*.,leftmargin=*]
\item Take a representative in the closed standard modular fundamental
region. If $T\ge7/5$, \ref{ari:C.8.13} applies. Otherwise
$|\Re\tau|\le1/2$ and $\sqrt3/2\le T\le7/5$. For fixed
$\Re\tau$, the function
\[
 |\chi|^2=\frac{(\Re\tau)^2+(T-1)^2}
                    {(\Re\tau)^2+(T+1)^2}
\]
has its maximum on this interval at an endpoint and increases
with $(\Re\tau)^2$. The two endpoint bounds are $r_0^2$
and $41/601<r_0^2$. Thus \ref{ari:C.8.14} covers the entire remaining
region. We conclude $|1-X|<1$ for every $t\ne1$.
Equations \eqref{ari:eq:cc-radial} now prove strict radial
monotonicity, including $t=0$ by continuity. They also exclude
both critical equations $W=0$ and $3Wp-2g_2=0$ on the
ordinary base. The elliptic points must be checked on their
root charts: at $\rho$, the nonzero coarse derivative already
given rules out a critical point; at $i$, $W=p=0$ but $3Wp-2g_2=-2g_2\ne0$, and C.3 gives
the strict maximum. At $\rho$, $W=g_2=0$ and $p\ne0$,
so both displayed factors vanish. This proves their exact zero
sets, including the elliptic orbits. In particular the vanishing of both factors at
$\rho$ in the $\tau$ coordinate is a ramification effect,
not a saddle on the coarse base. More explicitly, on the root
chart at $\rho$,
\[
 L=L(0)+\frac7{192}\Re\left(\frac{ip_0^2}{\pi^3}\delta^3\right)
                   +O(|\delta|^4),
\]
with $p_0\ne0$, whereas on the chart at $i$,
$L=\log\sqrt2-\pi|h|^2/2+O(|h|^4)$.
These expansions give the asserted root-chart critical types.

\item At the cusp, the product formula and \eqref{ari:eq:cc-EZ} give,
uniformly in $\Re\tau$,
\[
 L=-\frac{\pi T}{6}+\log(2\sqrt3)-
                      \frac{\pi\beta_{\rm cusp}^2}{T}+O(e^{-2\pi T}),
 \qquad |t|=\frac{e^{2\pi T}}{1728}(1+o(1)).
\]
Hence $L\to-\infty$ uniformly as $|t|\to\infty$.
Continuity at $t=1$ supplies the other endpoint of each ray.
The intermediate value theorem and strict monotonicity give
one and only one radius $r_c(\theta)>0$ at each level $c$.
Properness and strict monotonicity imply continuity in $\theta$;
the nonvanishing radial derivative gives the stated local
regularity by the implicit-function theorem. Thus
$\theta\mapsto1+r_c(\theta)e^{i\theta}$ is a continuous
injective parametrization of a Jordan curve and has exactly
the asserted complementary regions.

\item The proposed route in the stated conventions that posits one saddle on
$(0,1)$ is therefore false. The corrected statement is the
strict radial foliation above. It supplies, in addition to
the separate duplication proof, a global proof of Theorem~\ref{fam:src:C.2.3} and
resolves the critical-point and global-level-set question
posed for this normalized norm function.
\end{enumerate}
\end{proof}


